% Options for packages loaded elsewhere
\PassOptionsToPackage{unicode}{hyperref}
\PassOptionsToPackage{hyphens}{url}
\documentclass[
]{article}
\usepackage{xcolor}
\usepackage[margin=0.9in]{geometry}
\usepackage{amsmath,amssymb}
\setcounter{secnumdepth}{-\maxdimen} % remove section numbering
\usepackage{iftex}
\ifPDFTeX
  \usepackage[T1]{fontenc}
  \usepackage[utf8]{inputenc}
  \usepackage{textcomp} % provide euro and other symbols
\else % if luatex or xetex
  \usepackage{unicode-math} % this also loads fontspec
  \defaultfontfeatures{Scale=MatchLowercase}
  \defaultfontfeatures[\rmfamily]{Ligatures=TeX,Scale=1}
\fi
\usepackage{lmodern}
\ifPDFTeX\else
  % xetex/luatex font selection
\fi
% Use upquote if available, for straight quotes in verbatim environments
\IfFileExists{upquote.sty}{\usepackage{upquote}}{}
\IfFileExists{microtype.sty}{% use microtype if available
  \usepackage[]{microtype}
  \UseMicrotypeSet[protrusion]{basicmath} % disable protrusion for tt fonts
}{}
\makeatletter
\@ifundefined{KOMAClassName}{% if non-KOMA class
  \IfFileExists{parskip.sty}{%
    \usepackage{parskip}
  }{% else
    \setlength{\parindent}{0pt}
    \setlength{\parskip}{6pt plus 2pt minus 1pt}}
}{% if KOMA class
  \KOMAoptions{parskip=half}}
\makeatother
\usepackage{color}
\usepackage{fancyvrb}
\newcommand{\VerbBar}{|}
\newcommand{\VERB}{\Verb[commandchars=\\\{\}]}
\DefineVerbatimEnvironment{Highlighting}{Verbatim}{commandchars=\\\{\}}
% Add ',fontsize=\small' for more characters per line
\newenvironment{Shaded}{}{}
\newcommand{\AlertTok}[1]{\textcolor[rgb]{1.00,0.00,0.00}{\textbf{#1}}}
\newcommand{\AnnotationTok}[1]{\textcolor[rgb]{0.38,0.63,0.69}{\textbf{\textit{#1}}}}
\newcommand{\AttributeTok}[1]{\textcolor[rgb]{0.49,0.56,0.16}{#1}}
\newcommand{\BaseNTok}[1]{\textcolor[rgb]{0.25,0.63,0.44}{#1}}
\newcommand{\BuiltInTok}[1]{\textcolor[rgb]{0.00,0.50,0.00}{#1}}
\newcommand{\CharTok}[1]{\textcolor[rgb]{0.25,0.44,0.63}{#1}}
\newcommand{\CommentTok}[1]{\textcolor[rgb]{0.38,0.63,0.69}{\textit{#1}}}
\newcommand{\CommentVarTok}[1]{\textcolor[rgb]{0.38,0.63,0.69}{\textbf{\textit{#1}}}}
\newcommand{\ConstantTok}[1]{\textcolor[rgb]{0.53,0.00,0.00}{#1}}
\newcommand{\ControlFlowTok}[1]{\textcolor[rgb]{0.00,0.44,0.13}{\textbf{#1}}}
\newcommand{\DataTypeTok}[1]{\textcolor[rgb]{0.56,0.13,0.00}{#1}}
\newcommand{\DecValTok}[1]{\textcolor[rgb]{0.25,0.63,0.44}{#1}}
\newcommand{\DocumentationTok}[1]{\textcolor[rgb]{0.73,0.13,0.13}{\textit{#1}}}
\newcommand{\ErrorTok}[1]{\textcolor[rgb]{1.00,0.00,0.00}{\textbf{#1}}}
\newcommand{\ExtensionTok}[1]{#1}
\newcommand{\FloatTok}[1]{\textcolor[rgb]{0.25,0.63,0.44}{#1}}
\newcommand{\FunctionTok}[1]{\textcolor[rgb]{0.02,0.16,0.49}{#1}}
\newcommand{\ImportTok}[1]{\textcolor[rgb]{0.00,0.50,0.00}{\textbf{#1}}}
\newcommand{\InformationTok}[1]{\textcolor[rgb]{0.38,0.63,0.69}{\textbf{\textit{#1}}}}
\newcommand{\KeywordTok}[1]{\textcolor[rgb]{0.00,0.44,0.13}{\textbf{#1}}}
\newcommand{\NormalTok}[1]{#1}
\newcommand{\OperatorTok}[1]{\textcolor[rgb]{0.40,0.40,0.40}{#1}}
\newcommand{\OtherTok}[1]{\textcolor[rgb]{0.00,0.44,0.13}{#1}}
\newcommand{\PreprocessorTok}[1]{\textcolor[rgb]{0.74,0.48,0.00}{#1}}
\newcommand{\RegionMarkerTok}[1]{#1}
\newcommand{\SpecialCharTok}[1]{\textcolor[rgb]{0.25,0.44,0.63}{#1}}
\newcommand{\SpecialStringTok}[1]{\textcolor[rgb]{0.73,0.40,0.53}{#1}}
\newcommand{\StringTok}[1]{\textcolor[rgb]{0.25,0.44,0.63}{#1}}
\newcommand{\VariableTok}[1]{\textcolor[rgb]{0.10,0.09,0.49}{#1}}
\newcommand{\VerbatimStringTok}[1]{\textcolor[rgb]{0.25,0.44,0.63}{#1}}
\newcommand{\WarningTok}[1]{\textcolor[rgb]{0.38,0.63,0.69}{\textbf{\textit{#1}}}}
\usepackage{longtable,booktabs,array}
\newcounter{none} % for unnumbered tables
\usepackage{calc} % for calculating minipage widths
% Correct order of tables after \paragraph or \subparagraph
\usepackage{etoolbox}
\makeatletter
\patchcmd\longtable{\par}{\if@noskipsec\mbox{}\fi\par}{}{}
\makeatother
% Allow footnotes in longtable head/foot
\IfFileExists{footnotehyper.sty}{\usepackage{footnotehyper}}{\usepackage{footnote}}
\makesavenoteenv{longtable}
\setlength{\emergencystretch}{3em} % prevent overfull lines
\providecommand{\tightlist}{%
  \setlength{\itemsep}{0pt}\setlength{\parskip}{0pt}}
\usepackage[]{natbib}
\bibliographystyle{plainnat}
\usepackage{bookmark}
\IfFileExists{xurl.sty}{\usepackage{xurl}}{} % add URL line breaks if available
\urlstyle{same}
\hypersetup{
  hidelinks,
  pdfcreator={LaTeX via pandoc}}

\author{}
\date{}

\begin{document}

\section{From Agents to Institutions: A Field Study of Institutional Control in an AI-Staffed Prediction-Market Desk}\label{from-agents-to-institutions-a-field-study-of-institutional-control-in-an-ai-staffed-prediction-market-desk}

\textbf{Paper type:} Technical Report / Empirical Systems Study

\textbf{Author:} Wes Zheng

\textbf{Project repository:} https://github.com/wes-zheng/ai\_institutions

\subsection{Abstract}\label{abstract}

Modern agent systems increasingly support tool use, handoffs, memory, tracing, durable execution, guardrails, and multi-agent orchestration. These primitives improve agent capability, but they do not by themselves answer a different question: can AI workers produce accountable, authorized, reviewable, stoppable, and improvable work under operational pressure?

This paper reports a field study of an AI-staffed prediction-market desk under human-owner governance boundaries. The desk is used as a high-skill stress test rather than as a trading contribution. The domain forces the organization to handle uncertain evidence, changing external state, risk limits, no-action decisions, execution authority, delayed outcome truth, and post-outcome learning.

We use \textbf{institutional control} to mean the durable mechanisms by which an AI organization assigns ownership, constrains authority, records evidence, verifies outputs, validates closure, handles no-action, reconciles outcomes, and mutates doctrine. The central empirical finding is that many failures that appear to be agent failures are better diagnosed as organizational failures: missing ownership, broad intent not compiled into executable work, tool access mistaken for authority, plausible artifacts accepted without verifier state, stale messages treated as current work, completion confused with closure, and learning trapped in chat or tickets rather than becoming durable doctrine.

A central mechanism observed in the study is that the institutional layer did not appear fully formed: recurring fixes to repeated failures were promoted into durable role, work-record, message, verifier, replay, tool, and playbook artifacts by the manager employee, inside a human-owner governance boundary. We contribute a framework for AI organizations as a level above agents and workflows, bounded evidence from one human-owned deployment, and organization-level metrics for evaluating AI labor systems. The study does not claim causal certainty or deployment-wide performance improvement; it shows how institutional infrastructure can make AI labor more auditable, controllable, and better structured for governed improvement under operational pressure.

\subsection{1. Introduction}\label{introduction}

Agentic AI systems are usually evaluated through the behavior of a model, an agent, a workflow, or a team of agents. Those levels matter. A useful worker must be able to reason, use tools, plan, interact with external systems, and hand off work. But human labor is not performed by isolated intelligence. It is performed inside institutions that define who owns work, who has authority, what evidence is accepted, when work may stop, who verifies it, how closure is recorded, and how the organization changes after failure.

This paper argues that the path from agents to AI labor runs through institutions.

The distinction is practical. A capable agent can produce a plausible analysis while no one owns the next step. A tool-using agent can see market or account information without being authorized to act on it. A reviewer agent can provide helpful commentary without making an acceptance decision. A workflow can finish while the governing record remains unclosed. A memory system can store lessons while the next employee still acts from stale doctrine. These are not just prompt-engineering problems. They are organizational-control problems.

We study those problems through an AI-staffed prediction-market desk. The deployment was AI-staffed for daily operational labor and human-owned for goals, risk boundaries, and outer governance. It was composed of AI employees with durable roles, messages, work records, playbooks, review gates, execution boundaries, manager verification, replay records, and doctrine changes. The empirical domain is deliberately demanding. Weather and climate prediction-market work involves uncertainty modeling, source freshness, market interpretation, risk review, execution discipline, no-action judgment, and delayed settlement or outcome reconciliation. It also leaves artifacts when the organization fails.

``Manager employee'' refers to an AI employee role with authority to close, reroute, and perform scoped institutional mutations inside the human-owner governance boundary. The human owner retained goals, risk boundaries, and outer governance authority.

The paper is not about a trading strategy. It is about what had to be built around AI workers before their work became accountable, reviewable, stoppable, and improvable.

\subsubsection{1.1 From Agent Capability To Auditable Labor Control}\label{from-agent-capability-to-auditable-labor-control}

Agent-centered systems ask whether an AI can complete a task. Organization-centered systems ask whether the work can be owned, authorized, verified, stopped, closed, replayed, and adapted.

{\def\LTcaptype{none} % do not increment counter
\begin{longtable}[]{@{}
  >{\raggedright\arraybackslash}p{(\linewidth - 2\tabcolsep) * \real{0.5000}}
  >{\raggedright\arraybackslash}p{(\linewidth - 2\tabcolsep) * \real{0.5000}}@{}}
\toprule\noalign{}
\begin{minipage}[b]{\linewidth}\raggedright
Agent-centric question
\end{minipage} & \begin{minipage}[b]{\linewidth}\raggedright
Organization-centric question
\end{minipage} \\
\midrule\noalign{}
\endhead
\bottomrule\noalign{}
\endlastfoot
Can the agent act? & Who owns the work now? \\
Can the agent call tools? & What does tool evidence authorize? \\
Can agents hand off? & Did accountability transfer to the right owner? \\
Can a workflow finish? & Was closure valid and accepted? \\
Can memory persist? & Did learning become durable doctrine? \\
Can a reviewer critique output? & Did a verifier make a stateful accept/repair/reroute/decline decision? \\
\end{longtable}
}

The empirical system started with familiar agent ingredients: specialized roles, tools, work tracking, messages, and instructions. Over time, failures forced a more institutional design. Broad tasks became child work records. Personas became role contracts. Comments became message-trigger rules. Critique became verifier state. Completion became manager closure. Retros became replay and doctrine-mutation paths. The workboard became a state machine that compiled broad intent into live umbrellas, retro umbrellas, child paths, activations, gates, and replay obligations. Work records became the source of truth rather than conversation. Playbooks became operating doctrine rather than optional documentation. No-action became a valid output. Replay became the bridge from outcome truth to organizational learning.

\subsubsection{1.2 Thesis}\label{thesis}

The thesis is:

\begin{quote}
This study suggests that high-skill AI labor takes more than agent capability: it needs institutional-control mechanisms that make ownership, authority, evidence, verification, no-action, closure, replay, and governed adaptation auditable and controllable.
\end{quote}

This thesis does not imply that models, tools, memory, or workflow graphs are unimportant. It implies that they are insufficient units of analysis for high-skill labor. The same model and tool stack can behave very differently depending on the institution around it.

The organizing frame is four-part: institutional control converts agent outputs into auditable labor by governing state, authority, communication, and learning.

{\def\LTcaptype{none} % do not increment counter
\begin{longtable}[]{@{}
  >{\raggedright\arraybackslash}p{(\linewidth - 4\tabcolsep) * \real{0.3333}}
  >{\raggedright\arraybackslash}p{(\linewidth - 4\tabcolsep) * \real{0.3333}}
  >{\raggedright\arraybackslash}p{(\linewidth - 4\tabcolsep) * \real{0.3333}}@{}}
\toprule\noalign{}
\begin{minipage}[b]{\linewidth}\raggedright
Pillar
\end{minipage} & \begin{minipage}[b]{\linewidth}\raggedright
Main mechanisms
\end{minipage} & \begin{minipage}[b]{\linewidth}\raggedright
What failure it exposes or controls
\end{minipage} \\
\midrule\noalign{}
\endhead
\bottomrule\noalign{}
\endlastfoot
State control & work records, workboard, closure & chat-as-state and fake completion \\
Authority control & role contracts, verifier gates, tool boundaries & capability mistaken for permission \\
Communication control & messages, freshness checks, no-reply & false motion and stale triggers \\
Learning control & replay, doctrine mutation, case-only classification & lessons trapped in memory or tickets \\
\end{longtable}
}

\subsubsection{1.3 Research Questions}\label{research-questions}

This study asks five questions:

\textbf{RQ1. Organizational primitives:} What primitives are needed to turn agent-like workers into AI employees inside a durable organization?

\textbf{RQ2. Auditability and controllability:} Which organizational primitives make ownership, authority, evidence, closure, and learning failures visible, bounded, or controllable?

\textbf{RQ3. Decision work:} Can an AI organization produce verifier-reconstructible technical artifacts rather than plausible prose?

\textbf{RQ4. Learning:} How does delayed outcome truth become durable organizational change rather than one-off correction?

\textbf{RQ5. Generalization:} Which design principles transfer beyond prediction-market trading to other high-skill technical organizations?

\subsubsection{1.4 Contributions}\label{contributions}

This paper makes three contributions.

\begin{enumerate}
\def\labelenumi{\arabic{enumi}.}
\tightlist
\item
  \textbf{Conceptual contribution:} It defines AI organization and institutional control as an evaluation level above models, agents, workflows, and teams.
\item
  \textbf{Empirical contribution:} It reports field evidence from a human-owned, AI-staffed high-skill operational desk showing ownership, authority, evidence, workboard-state, topology, closure, no-action, replay, and mutation failures.
\item
  \textbf{Evaluation contribution:} It proposes organization-level metrics and sanitized artifact field sets for assessing AI labor systems without reducing evaluation to task completion or model capability.
\end{enumerate}

\subsection{2. Related Work And Positioning}\label{related-work-and-positioning}

The public baseline as of April 2026 is strong. Agent frameworks have matured rapidly. Tool protocols have become standard infrastructure. Enterprise agent platforms increasingly expose governed deployment primitives. Benchmarks increasingly evaluate long-horizon web, computer-use, software-engineering, workplace, and tool-use tasks. Multi-agent systems have explored role decomposition, group chat, social simulation, and standard operating procedures.

This paper does not claim those directions are wrong. They provide increasingly powerful execution, communication, governance, and deployment primitives, but public systems and benchmarks still rarely treat organization-level labor state as the primary unit of empirical evaluation.

To our knowledge, this is an early field study to treat an AI-staffed organization's durable labor artifacts as the empirical unit of analysis, with bounded evidence for ownership, authority, no-action, manager closure, replay, and doctrine mutation.

These systems improve execution. This paper studies institutional acceptance, authority, closure, no-action, replay, and doctrine mutation.

\subsubsection{2.1 Agent SDKs And Execution Infrastructure}\label{agent-sdks-and-execution-infrastructure}

OpenAI's Agents SDK and related agent tooling provide primitives for agents, tools, handoffs, guardrails, tracing, sandboxed execution, skills, and durable execution. OpenAI's April 2026 Agents SDK update describes a model-native harness, tool use through MCP, sandboxed execution, progressive disclosure through skills, and checkpoint-like restoration of agent state after environment loss \citep{openai2026agentsSdkEvolution}. The Agents SDK documentation describes tracing of model calls, tool calls, handoffs, and guardrails \citep{openai2026agentsSdkTracing}, including input, output, and tool guardrails for custom function-tool invocations \citep{openai2026agentsSdkGuardrails}.

These are valuable execution and validation primitives: they help developers run, observe, interrupt, and constrain agent workflows. But they do not by themselves define whether a worker has authority to close a work item, whether a no-action outcome is valid, whether an old message is superseded, whether a verifier must accept or reroute an artifact, or where reusable learning must land.

\subsubsection{2.2 Tool Protocols And MCP}\label{tool-protocols-and-mcp}

The Model Context Protocol, introduced by Anthropic as an open standard for connecting AI assistants to external data sources and tools, standardizes tool and context integration for AI applications \citep{anthropic2024modelContextProtocol}. The current MCP specification defines a JSON-RPC protocol with hosts, clients, servers, resources, prompts, tools, sampling, progress, cancellation, logging, and security considerations \citep{mcp2025specification}. Related 2025-11-25 specification material emphasizes schema conformance and protocol metadata \citep{mcp2025basicSpecification}. Anthropic later donated MCP to the Agentic AI Foundation, where it joined other agentic infrastructure projects under neutral stewardship \citep{anthropic2025mcpAaifDonation}. Recent MCP production-pattern work, published as a preprint, argues that production deployments still need additional mechanisms such as identity propagation, adaptive tool/timeout budgeting, and structured error semantics \citep{srinivasan2026bridgingProtocolProduction}. This report addresses a separate layer above tool integration: institutional authority, work ownership, verifier state, no-action, replay, and doctrine mutation.

MCP is a major infrastructure layer, but a protocol for tool access is not an organization. It does not, by itself, decide that an execution employee may only act inside a risk-authorized envelope, that a reviewer can decline but not execute, which role may mutate durable doctrine inside a human-owner governance boundary, or that a message is stale because a governing work record has moved on. This paper positions organization-protocol semantics above MCP: employee identity, role authority, work state, doctrine, message lifecycle, verification, replay, and institutional mutation.

\subsubsection{2.3 Enterprise Agent Platforms And Agent-To-Agent Protocols}\label{enterprise-agent-platforms-and-agent-to-agent-protocols}

Recent enterprise agent platforms increasingly expose organizational controls such as sharing, permissions, approval gates, audit logs, observability, agent identity, and lifecycle management \citep{openai2026workspaceAgents, googleCloud2026geminiEnterpriseAgents, googleCloud2026agentIdentity, microsoft2026copilotStudioAgentReadiness, salesforce2026agentforcePlatform}. A2A standardizes inter-agent communication, discovery, and task coordination across heterogeneous agents and platforms \citep{google2025agent2agent, linuxFoundation2025agent2agentProject}. These developments strengthen the premise that agentic work is becoming organizational. This report studies a narrower empirical layer: whether durable labor artifacts make ownership, authority, no-action, closure, replay, and doctrine mutation auditable in a live AI-staffed organization.

\subsubsection{2.4 Durable Workflow And Multi-Agent Frameworks}\label{durable-workflow-and-multi-agent-frameworks}

LangGraph provides durable execution and persistence for long-running workflows, including human-in-the-loop and interruption recovery \citep{langchain2026langgraphDurableExecution}. CrewAI combines agents, crews, flows, tasks, guardrails, callbacks, triggers, persistence, and human-in-the-loop patterns \citep{crewai2026documentation}. AutoGen and related systems explore multi-agent conversations, teams, handoffs, and event-driven workflows \citep{wu2023autogen}. MetaGPT encodes standard operating procedures into prompt sequences and role-based agent collaboration for software work \citep{hong2023metagpt}.

These systems move beyond a single agent, but the primary abstraction remains workflow, team, or collaboration. An AI organization adds stronger semantics: not only that an agent handed off, but whether accountability moved to the correct owner; not only that a workflow resumed, but whether the governing work record is truthful; not only that a role exists, but whether the role owns or is forbidden from a decision; not only that an SOP was available, but whether reusable learning changed durable doctrine without losing prior constraints.

\subsubsection{2.5 Social Simulation And Agent Memory}\label{social-simulation-and-agent-memory}

Generative Agents demonstrated believable social behavior by storing natural-language experiences, synthesizing reflections, and retrieving memories to plan behavior in a simulated town \citep{park2023generativeAgents}. That line of work is important for humanlike interaction and social simulation. This paper takes a different stance: humanlike behavior is not the goal. Accountable institutional labor is the goal. A worker can be socially plausible while still failing to own, authorize, verify, close, or learn correctly.

\subsubsection{2.6 Agent Benchmarks And Evaluation}\label{agent-benchmarks-and-evaluation}

Agent benchmarks increasingly evaluate more realistic tasks. WebArena provides realistic web environments and long-horizon web tasks, reporting large gaps between agents and humans in early systems \citep{zhou2023webarena}. OSWorld evaluates computer-use agents in real desktop environments \citep{xie2024osworld}. SWE-bench evaluates whether language models can resolve real GitHub issues \citep{jimenez2023swebench}. Tau-bench evaluates tool-agent-user interaction in domains with policies and tools, emphasizing reliability over repeated trials \citep{yao2024taubench}. AgentBench evaluates LLMs as agents across multiple interactive environments \citep{liu2023agentbench}, and TheAgentCompany moves closer to workplace labor by evaluating agents in a simulated software-company environment \citep{xu2024theagentcompany}. WorkArena and WorkArena++ further move agent evaluation toward enterprise knowledge work, but still primarily evaluate task completion rather than institutional state, authority, closure, and learning destinations \citep{drouin2024workarena, boisvert2024workarenaPlus}. A survey of LLM-agent evaluation, submitted in 2025 and revised in 2026, identifies planning, tool use, memory, application benchmarks, generalist agents, cost-efficiency, safety, robustness, and scalable fine-grained evaluation as major themes \citep{yehudai2025surveyAgents}.

These benchmarks strengthen the shift from static task evaluation toward agentic and workplace task performance. This report asks a different question: whether AI labor is owned, authorized, closed, stopped, replayed, and mutated as institutional state. An agent can pass a task benchmark while still failing organizationally; conversely, an organization can correctly stop a task that a benchmark might treat as incomplete.

\subsubsection{2.7 Business Process Management And Process Mining}\label{business-process-management-and-process-mining}

Business process management and process-mining systems model tasks, events, handoffs, conformance, and process improvement \citep{dumas2018fundamentalsBpm, vanDerAalst2016processMining}. This report is related, but the control problem differs because probabilistic AI workers require explicit role authority, evidence demotion, message freshness, valid no-action, verifier acceptance, and replay-to-doctrine mutation around model calls.

\subsubsection{2.8 Institutions, Norms, And Agent Governance}\label{institutions-norms-and-agent-governance}

The paper also sits in a longer multi-agent systems tradition. Normative multi-agent systems study how norms, roles, enforcement, and compliance shape autonomous-agent interaction \citep{boella2006normativeMultiagentSystems}. Electronic institutions make conventions explicit for open multi-agent systems and use institutional artifacts to shape participation and compliance \citep{esteva2004electronicInstitutions}.

Recent LLM-agent governance preprints and specifications are closer still. Institutional AI reframes multi-agent alignment from preference engineering in agent-space to mechanism design in institution-space, using governance graphs, public manifests, legal states, transitions, sanctions, restorative paths, and audit logs \citep{bracaleSyrnikov2026institutionalCournot, pierucci2026institutionalDistributional}. Agent Control Protocol frames agent action as a temporal admission-control problem between agent intent and system state mutation, emphasizing identity, capability scope, delegation, risk, execution history, and auditability \citep{fernandez2026agentControlProtocol, acp2026architecture}.

This paper is complementary. It does not propose a formal admission-control protocol or a governance graph. It brings the institutional lens to empirical AI labor: work records, messages, role authority, verifier topology, no-action, closure, replay, and doctrine mutation in a live operational desk.

\subsubsection{2.9 Organizational Learning And High-Reliability Work}\label{organizational-learning-and-high-reliability-work}

Organizational-learning and high-reliability work study how institutions convert operational experience into procedures, roles, checks, and routines \citep{argyrisSchon1978organizationalLearning, weickSutcliffe2015managingUnexpected}. This report adapts that lens to AI labor: replay records, role/playbook versions, doctrine patches, verifier gates, tool-parity repairs, and activation messages are treated as AI-native organizational-learning artifacts.

We use organizational-learning and high-reliability work as positioning lenses, not as evidence that this deployment achieved high-reliability-organization status. The mechanism claim is narrower: durable learning must land in institutional artifacts rather than remaining in private memory, one-off reflection, chat, or ticket commentary.

\subsubsection{2.10 Positioning Summary}\label{positioning-summary}

{\def\LTcaptype{none} % do not increment counter
\begin{longtable}[]{@{}
  >{\raggedright\arraybackslash}p{(\linewidth - 4\tabcolsep) * \real{0.3333}}
  >{\raggedright\arraybackslash}p{(\linewidth - 4\tabcolsep) * \real{0.3333}}
  >{\raggedright\arraybackslash}p{(\linewidth - 4\tabcolsep) * \real{0.3333}}@{}}
\toprule\noalign{}
\begin{minipage}[b]{\linewidth}\raggedright
Public baseline
\end{minipage} & \begin{minipage}[b]{\linewidth}\raggedright
What it provides
\end{minipage} & \begin{minipage}[b]{\linewidth}\raggedright
Organization-level gap addressed here
\end{minipage} \\
\midrule\noalign{}
\endhead
\bottomrule\noalign{}
\endlastfoot
Agent SDKs & tools, handoffs, guardrails, tracing, sandboxing & durable authority, closure, doctrine, no-action, replay \\
MCP & standardized tool/resource/prompt access & employee identity, work ownership, role authority, verifier state \\
A2A / agent-to-agent protocols & agent discovery, communication, and task coordination & institutional ownership, authority, closure, replay, and mutation semantics \\
Enterprise agent platforms & sharing, permissions, approvals, governance, observability, agent scaling & durable empirical labor artifacts for ownership, no-action, closure, replay, doctrine mutation \\
Workflow graphs & persistence, routing, interrupts & organizational truth, manager finality, institutional mutation \\
BPM and process mining & task/event models, handoffs, conformance, process improvement & AI-specific authority, evidence demotion, no-action, and replay-to-doctrine mutation \\
Multi-agent teams & collaboration and role decomposition & authority contracts, stateful handoffs, accepted closure \\
Agent benchmarks & task success in environments & organization-level auditability and learning \\
Workplace-agent benchmarks & workplace-like task performance & organization-level state and learning evidence \\
Institutional AI / ACP & formal governance graphs, admission control, auditability & empirical labor records, closure, no-action, replay, doctrine mutation \\
Organizational learning and high-reliability work & experience-to-routine learning and resilience practices & machine-executable replay, versioned doctrine, and activation-gated mutation \\
\end{longtable}
}

The paper's claim is not that agent systems lack infrastructure. The claim is that high-skill AI labor may need an additional institutional layer.

\subsection{3. Framework: AI Organization And Institutional Control}\label{framework-ai-organization-and-institutional-control}

\subsubsection{3.1 Level Hierarchy}\label{level-hierarchy}

The proposed hierarchy is:

{\def\LTcaptype{none} % do not increment counter
\begin{longtable}[]{@{}
  >{\raggedright\arraybackslash}p{(\linewidth - 4\tabcolsep) * \real{0.3333}}
  >{\raggedright\arraybackslash}p{(\linewidth - 4\tabcolsep) * \real{0.3333}}
  >{\raggedright\arraybackslash}p{(\linewidth - 4\tabcolsep) * \real{0.3333}}@{}}
\toprule\noalign{}
\begin{minipage}[b]{\linewidth}\raggedright
Level
\end{minipage} & \begin{minipage}[b]{\linewidth}\raggedright
Primary question
\end{minipage} & \begin{minipage}[b]{\linewidth}\raggedright
Limitation for auditable labor control
\end{minipage} \\
\midrule\noalign{}
\endhead
\bottomrule\noalign{}
\endlastfoot
Model & Can it reason or generate? & no durable ownership or authority \\
Agent & Can it act with tools? & weak institutional accountability \\
Workflow & Can steps be routed? & weak doctrine, authority, and learning semantics \\
Team & Can specialists coordinate? & conversation can masquerade as state \\
Organization & Can work be owned, authorized, verified, stopped, replayed, and adapted? & target level \\
\end{longtable}
}

Figure 1. From Models To AI Organizations.

\begin{Shaded}
\begin{Highlighting}[]
\NormalTok{model}
\NormalTok{  {-}\textgreater{} agent}
\NormalTok{  {-}\textgreater{} workflow / team}
\NormalTok{  {-}\textgreater{} organization}

\NormalTok{unit of control:}
\NormalTok{generation {-}\textgreater{} tool action {-}\textgreater{} routed steps {-}\textgreater{} institutional labor}
\end{Highlighting}
\end{Shaded}

Takeaway: the paper studies the organizational layer where ownership, authority, closure, replay, and mutation become evaluable.

An organization is not just ``many agents.'' It is a durable control system around workers.

\subsubsection{3.2 Definition}\label{definition}

An \textbf{AI organization} is a durable socio-technical control system in which AI employees perform bounded work under explicit role contracts, authority limits, governing work records, reusable doctrine, verifier gates, and outcome-driven learning loops.

``AI employee'' is used here as an operational term for a persistent institutional principal with role, authority, inbox, tools, and history. It is not a legal, moral, or anthropomorphic claim.

This definition deliberately separates worker capability from organizational controllability. The organization is responsible for converting probabilistic model outputs into accountable labor artifacts.

\subsubsection{3.3 Institutional Control}\label{institutional-control}

\textbf{Institutional control} is the durable mechanism set by which an AI organization assigns ownership, constrains authority, records evidence, verifies outputs, validates closure, handles no-action, reconciles outcomes, and mutates doctrine. In this operational sense, it rests on six negations:

\begin{enumerate}
\def\labelenumi{\arabic{enumi}.}
\tightlist
\item
  The model is not the system of record.
\item
  Conversation is not durable organizational state.
\item
  Tool access is not authority.
\item
  Completion is not self-declared.
\item
  Memory is not doctrine.
\item
  Reflection is not learning until it changes the right durable layer or is explicitly classified as case-only.
\end{enumerate}

Persistence is treated as substrate, not as a standalone novelty. The relevant institutional state is employee identity, role text, playbooks, inboxes, governing work records, tool authority, and replay obligations. Private memory may help continuity, but it is a hint, not the system of record. When private memory conflicts with the governing record, role text, or playbook doctrine, the durable institutional surface controls.

Institutional control is therefore less about making a single worker more autonomous and more about making a group of workers governable.

\subsubsection{3.4 Organizational Primitives}\label{organizational-primitives}

{\def\LTcaptype{none} % do not increment counter
\begin{longtable}[]{@{}
  >{\raggedright\arraybackslash}p{(\linewidth - 4\tabcolsep) * \real{0.3333}}
  >{\raggedright\arraybackslash}p{(\linewidth - 4\tabcolsep) * \real{0.3333}}
  >{\raggedright\arraybackslash}p{(\linewidth - 4\tabcolsep) * \real{0.3333}}@{}}
\toprule\noalign{}
\begin{minipage}[b]{\linewidth}\raggedright
Primitive
\end{minipage} & \begin{minipage}[b]{\linewidth}\raggedright
Function
\end{minipage} & \begin{minipage}[b]{\linewidth}\raggedright
Failure mode exposed or controlled
\end{minipage} \\
\midrule\noalign{}
\endhead
\bottomrule\noalign{}
\endlastfoot
Employee principal & Persistent worker identity, inbox, role, tools, status & stateless task execution \\
Role contract & Owned decisions, forbidden decisions, outputs, escalation & persona drift and unauthorized action \\
Governing work record & Case truth: owner, status, evidence, blocker, next action & chat-as-state and fake progress \\
Workboard state machine & Intent compiled into umbrellas, child records, gate states, owners, activations, closure, and replay & vague intent that never becomes executable labor \\
Work message & Trigger and handoff with exact ask and reply expectation & missed activation and acknowledgment loops \\
Playbook doctrine & Reusable workflow and evidence rules & repeated one-off correction \\
Evidence surface & Accepted source/tool outputs attached to records & unverifiable reasoning \\
Verifier gate & Accept, repair, reroute, decline, block, or approve & plausible but incomplete artifacts \\
No-action state & Valid stop, decline, stand-down, already-done, explicit zero & forced action and false completion \\
Replay record & Outcome reconciliation and learning classification & learning detached from truth \\
Institutional mutation & Changes to roles, playbooks, tools, staffing, topology, workboards, or messages & repeated failure risk, unactivated fixes, and learning trapped in local context \\
\end{longtable}
}

Figure 2. Institutional Control Stack.

\begin{Shaded}
\begin{Highlighting}[]
\NormalTok{human{-}owner boundary}
\NormalTok{  {-}\textgreater{} organization topology and manager finality}
\NormalTok{  {-}\textgreater{} employee principal + role contract}
\NormalTok{  {-}\textgreater{} playbook doctrine}
\NormalTok{  {-}\textgreater{} governing work record}
\NormalTok{  {-}\textgreater{} message trigger + freshness gate}
\NormalTok{  {-}\textgreater{} evidence packet + verifier gate}
\NormalTok{  {-}\textgreater{} closure / no{-}action / replay / mutation}
\end{Highlighting}
\end{Shaded}

Takeaway: the model invocation is only one worker cycle inside a larger control stack.

\subsubsection{3.5 Propositions}\label{propositions}

The empirical study is organized around five propositions.

\textbf{Proposition 1. Ownership:} AI work becomes more auditable and controllable when every substantive work item has a governing record with a current owner and exact next action.

\textbf{Proposition 2. Authority:} AI work is more controllable when role authority is separated from tool access and downstream action depends on explicit authorization.

\textbf{Proposition 3. Evidence:} AI decision artifacts become more reviewable when evidence, assumptions, and decision structure are stored outside private model context.

\textbf{Proposition 4. No-action:} AI organizations benefit from first-class states for repair, decline, blocked, stand-down, already-done, no-send, and explicit zero.

\textbf{Proposition 5. Replay:} Outcome reconciliation becomes more traceable when replay is linked to the role, playbook, evidence, and work-state artifacts active at the time of action.

\subsection{4. Empirical Setting And Study Design}\label{empirical-setting-and-study-design}

\subsubsection{4.1 Empirical Setting}\label{empirical-setting}

The empirical setting is an AI-staffed prediction-market desk working in weather and climate markets. The organization performed research, evidence collection, market-path selection, risk review, execution or no-action decisions, manager closure, and replay. It was AI-staffed for daily operational labor and human-owned for goals, risk boundaries, and outer governance.

The domain is useful because it stresses several properties common to high-skill technical work:

\begin{itemize}
\tightlist
\item
  external evidence changes over time;
\item
  intermediate artifacts must be verifiable by another role;
\item
  action can be unsupported even when analysis is plausible;
\item
  stopping can be correct;
\item
  delayed outcome truth can evaluate earlier decisions;
\item
  organizational mistakes leave durable traces in work records, messages, reviews, and replay.
\end{itemize}

The desk is not presented as evidence that AI systems should trade, nor as financial advice. It is a field setting where institutional failures are observable.

\subsubsection{4.1.1 Pre-Institutional Reference Condition}\label{pre-institutional-reference-condition}

The reference condition was not a separate model architecture or controlled benchmark baseline. It was the earlier operating mode of the same deployment: role-like agents using tools, work tracking, messages, and instructions before mature ownership, authority, closure, no-action, replay, and mutation semantics were introduced.

The study therefore compares observed failures in that pre-institutional reference condition with bounded mature-state audits after institutional-control mechanisms were introduced. It should not be read as a randomized comparison or as a benchmark-style model evaluation.

\subsubsection{4.2 Unit Of Analysis}\label{unit-of-analysis}

The primary empirical unit is a \textbf{market-path episode}: a bounded path from intake through research, review, action or no-action, closure, and replay where applicable.

Episodes may end in several valid states:

\begin{itemize}
\tightlist
\item
  no-send or no-action;
\item
  repair request;
\item
  risk decline;
\item
  execution stand-down;
\item
  executed action followed by replay;
\item
  manager closure;
\item
  durable doctrine or role/tool/topology mutation;
\item
  case-only learning classification.
\end{itemize}

This matters because the organization is not evaluated only on actions taken. It is evaluated on whether it reached the truthful state.

\subsubsection{4.3 Data Artifacts}\label{data-artifacts}

The empirical record consists of:

\begin{itemize}
\tightlist
\item
  governing work records;
\item
  task and handoff messages;
\item
  role contracts and role versions;
\item
  playbooks and playbook versions;
\item
  verifier decisions;
\item
  manager closure records;
\item
  replay and retro records;
\item
  owner-directed change records;
\item
  tool/profile/tool-set records;
\item
  realized outcome records where available.
\end{itemize}

This report uses sanitized extraction summaries, not raw internal records. A first evidence extraction pass deep-read selected work records and comment threads, performed mechanism keyword sweeps, and produced sanitized evidence packets. A message evidence pass sampled manager, execution-lead, and verifier inboxes for no-reply suppression, self-trigger discipline, stale/unauthorized reroute, and verifier decision patterns. A bounded sprint-window pass extracted denominator-labeled counts for selected Results metrics.

\subsubsection{4.4 Study Scope And Evidence Corpus}\label{study-scope-and-evidence-corpus}

This is a field study and mechanism paper. It does not claim general causal certainty. It identifies recurring organizational failure modes and the control mechanisms that made them observable, bounded, or repairable.

The deployment context is summarized below. The full corpus is not used as a denominator in this report; bounded counts are reported only for selected windows with enough sanitized evidence to code the relevant fields.

The evidence hierarchy is:

{\def\LTcaptype{none} % do not increment counter
\begin{longtable}[]{@{}
  >{\raggedright\arraybackslash}p{(\linewidth - 4\tabcolsep) * \real{0.3333}}
  >{\raggedright\arraybackslash}p{(\linewidth - 4\tabcolsep) * \real{0.3333}}
  >{\raggedright\arraybackslash}p{(\linewidth - 4\tabcolsep) * \real{0.3333}}@{}}
\toprule\noalign{}
\begin{minipage}[b]{\linewidth}\raggedright
Evidence class
\end{minipage} & \begin{minipage}[b]{\linewidth}\raggedright
What it supports
\end{minipage} & \begin{minipage}[b]{\linewidth}\raggedright
What it does not support
\end{minipage} \\
\midrule\noalign{}
\endhead
\bottomrule\noalign{}
\endlastfoot
S1 bounded audit & ownership, closure, no-action, replay visibility & global rates or correctness \\
P1/B1 spot-checks & selection-bias context & random-sample inference \\
M1 message packets & false-motion and freshness mechanisms & inbox-wide loop rate \\
Verifier comments & repair, approval, reroute, decline mechanisms & full verifier-quality rate \\
Case studies & mechanism tracing & recurrence reduction \\
\end{longtable}
}

{\def\LTcaptype{none} % do not increment counter
\begin{longtable}[]{@{}
  >{\raggedright\arraybackslash}p{(\linewidth - 2\tabcolsep) * \real{0.5000}}
  >{\raggedright\arraybackslash}p{(\linewidth - 2\tabcolsep) * \real{0.5000}}@{}}
\toprule\noalign{}
\begin{minipage}[b]{\linewidth}\raggedright
Deployment attribute
\end{minipage} & \begin{minipage}[b]{\linewidth}\raggedright
Value
\end{minipage} \\
\midrule\noalign{}
\endhead
\bottomrule\noalign{}
\endlastfoot
Full deployment period & Multi-week deployment; deep audit windows cover Apr 25-27, 2026, with exploratory spot-check records from earlier periods \\
AI employee roles & Multiple persistent AI employee principals across the role families below; full roster count is not reported \\
Active role families & Qualitatively includes management, research, verification, execution, replay, doctrine, and tooling functions \\
Work records in full corpus & More than 2,000 internal work records exist; not used as a corpus-wide denominator here \\
Messages in full corpus & Full count not estimated; \texttt{M1} requested 260 inbox rows for mechanism evidence \\
Playbook versions & Versioned doctrine exists; count-level estimate is not reported \\
Models used & LLM API models; exact provider/configuration is not disclosed in this report and not used as causal evidence \\
Human-owner interventions & Bounded S1 sample includes one owner directive; full-corpus intervention count is not estimated \\
\end{longtable}
}

The bounded empirical counts in this report come from selected windows and message samples.

The selected windows are bounded mature-state audits, not a corpus-wide performance estimate of the entire deployment.

{\def\LTcaptype{none} % do not increment counter
\begin{longtable}[]{@{}
  >{\raggedright\arraybackslash}p{(\linewidth - 6\tabcolsep) * \real{0.2308}}
  >{\raggedleft\arraybackslash}p{(\linewidth - 6\tabcolsep) * \real{0.3077}}
  >{\raggedright\arraybackslash}p{(\linewidth - 6\tabcolsep) * \real{0.2308}}
  >{\raggedright\arraybackslash}p{(\linewidth - 6\tabcolsep) * \real{0.2308}}@{}}
\toprule\noalign{}
\begin{minipage}[b]{\linewidth}\raggedright
Artifact type
\end{minipage} & \begin{minipage}[b]{\linewidth}\raggedleft
Count used in this report
\end{minipage} & \begin{minipage}[b]{\linewidth}\raggedright
Window / source
\end{minipage} & \begin{minipage}[b]{\linewidth}\raggedright
Used for
\end{minipage} \\
\midrule\noalign{}
\endhead
\bottomrule\noalign{}
\endlastfoot
Market-path live child records & 46 & S1, R1, V1 selected sprint windows, Apr 25-27 2026 & ownership, verifier, no-action, replay linkage \\
Work records in selected windows & 80 & live children, replay children, and live/retro umbrellas & state control, closure, replay, parent/child truth \\
S1 final closeout records & 20 & S1 final manager verification sample & runtime-contract compliance, closure validity, no-message closeout \\
S1 no-execution paths & 6 & S1 final live dispositions & no-action quality and no-replay rationale \\
S1 ordered/replay paths & 14 & S1 replay children & replay coverage and delayed outcome truth \\
Replay child records across S1/R1/V1 & 28 & selected sprint/retro windows & replay learning and no-replay separation \\
Message rows requested & 260 & manager, execution-lead, and verifier inbox sample & no-reply, stale/reroute, self-trigger, verifier-decision examples \\
Message evidence packets & 6 & high-signal message extraction & message-trigger discipline and loop suppression \\
S1 functional verifier approvals & at least 16 & S1 functional-verifier comments & verifier topology and approval behavior \\
S1 repair-required verifier decisions & at least 10 & S1 functional-verifier comments & evidence sufficiency and repair routing \\
S1 observed reroutes & 4 & S1 comments & stateful verifier/manager correction \\
S1 changed-envelope approvals & 3 & S1 comments & verifier decisions that changed downstream action \\
S1 human-owner directive & 1 & S1 umbrella open directive & human-owner boundary and setup role \\
Durable mutation artifacts & qualitative mechanism evidence only & selected evolution/tool records & supports mechanism description, not aggregate mutation-rate claims \\
P1 pre-institutional/reference spot-check & 30 & earlier work-record metadata sample & selected-window bias check and reference contrast \\
B1 broader non-S1/R1/V1 spot-check & 75 & shallow metadata sample across other periods/statuses & mechanism-presence check outside selected mature windows \\
\end{longtable}
}

The study does not isolate model-version effects. Because model identity and configuration are not used as causal evidence, all claims are limited to institutional auditability and mechanism tracing rather than model-independent performance improvement.

\subsubsection{4.5 Sample Windows}\label{sample-windows}

The Results section uses three named sprint windows:

\begin{itemize}
\tightlist
\item
  \textbf{S1:} Apr 26 20-location live sprint plus retro. This is the primary bounded sample: 20 live child records, 14 ordered/replay paths, 6 no-execution/no-replay paths, and 20 final manager-reviewed live dispositions.
\item
  \textbf{R1:} Apr 25 daily-high live sprint plus retro. This is replay and parent-closeout augmentation: 20 live child records and 9 replay children.
\item
  \textbf{V1:} Apr 27 mini sprint plus retro. This is a compact recency check: 6 live child records and 5 replay children.
\end{itemize}

The message evidence sample \texttt{M1} requested 260 inbox rows across manager, execution-lead, and verifier inboxes and produced 6 sanitized mechanism packets. It is not a full inbox census.

\subsubsection{4.6 Sampling Rationale}\label{sampling-rationale}

The selected windows were chosen because they represent mature-state operation after the relevant institutional-control mechanisms were active and because they contained complete live, closeout, and replay artifacts. They are not claimed to be representative of the full deployment. Their purpose is to test whether the mature control stack made ownership, authority, no-action, closure, and replay auditable in bounded operational episodes.

Episodes were excluded from bounded metrics when the work path lacked enough sanitized evidence to code owner, state, verifier decision, or replay status. This exclusion rule supports cleaner denominator-labeled audits, but it also means the sample should not be read as a full deployment rate estimate.

\subsubsection{4.7 Broader Spot-Check Against Selection Bias}\label{broader-spot-check-against-selection-bias}

To reduce the risk that the paper only reports clean mature-state windows, we added a shallow broader spot-check. \texttt{P1} was selected from the earliest available substantive weather/work-record search results in late-March and early-April records, excluding administrative-only records when obvious. \texttt{B1} was selected from non-S1/R1/V1 records across multiple earlier operational periods and statuses, including live child, replay, duplicate/canceled, done, todo, and backlog records where available. These are not random samples from a formal database export; they are availability-bias checks against the selected mature windows.

Records were coded from work-record metadata and description snippets only. A field counted as clear only when the record exposed the relevant owner, next action or blocker, closure state, verifier state, replay/no-replay state, or no-action state without needing private comment-thread context. These codes measure state visibility, not decision correctness.

Owner-present coding means the record exposed the accountable current owner for the next institutional step; it does not mean other employees were free to act on that work.

{\def\LTcaptype{none} % do not increment counter
\begin{longtable}[]{@{}
  >{\raggedright\arraybackslash}p{(\linewidth - 14\tabcolsep) * \real{0.0968}}
  >{\raggedleft\arraybackslash}p{(\linewidth - 14\tabcolsep) * \real{0.1290}}
  >{\raggedleft\arraybackslash}p{(\linewidth - 14\tabcolsep) * \real{0.1290}}
  >{\raggedleft\arraybackslash}p{(\linewidth - 14\tabcolsep) * \real{0.1290}}
  >{\raggedleft\arraybackslash}p{(\linewidth - 14\tabcolsep) * \real{0.1290}}
  >{\raggedleft\arraybackslash}p{(\linewidth - 14\tabcolsep) * \real{0.1290}}
  >{\raggedleft\arraybackslash}p{(\linewidth - 14\tabcolsep) * \real{0.1290}}
  >{\raggedleft\arraybackslash}p{(\linewidth - 14\tabcolsep) * \real{0.1290}}@{}}
\toprule\noalign{}
\begin{minipage}[b]{\linewidth}\raggedright
Sample
\end{minipage} & \begin{minipage}[b]{\linewidth}\raggedleft
Records
\end{minipage} & \begin{minipage}[b]{\linewidth}\raggedleft
Owner present
\end{minipage} & \begin{minipage}[b]{\linewidth}\raggedleft
Next clear
\end{minipage} & \begin{minipage}[b]{\linewidth}\raggedleft
Closure clear
\end{minipage} & \begin{minipage}[b]{\linewidth}\raggedleft
Verifier clear
\end{minipage} & \begin{minipage}[b]{\linewidth}\raggedleft
Replay clear
\end{minipage} & \begin{minipage}[b]{\linewidth}\raggedleft
No-action clear
\end{minipage} \\
\midrule\noalign{}
\endhead
\bottomrule\noalign{}
\endlastfoot
\texttt{P1} reference spot-check & 30 & 21/30 & 18/30 & 22/30 & 10/18; N/A 12 & 8/14; N/A 16 & 6/13; N/A 17 \\
\texttt{S1} mature deep audit & 20 & 20/20 & 20/20 & 20/20 & 20/20 & 20/20 & 6/6; N/A 14 \\
\texttt{B1} broader spot-check & 75 & 62/75 & 51/75 & 58/75 & 34/55; N/A 20 & 30/46; N/A 29 & 12/26; N/A 49 \\
\end{longtable}
}

For S1, verifier clear refers to manager-reviewed final dispositions; replay clear means each path had replay or explicit no-replay classification. The spot-check supports two limited claims. First, S1 is cleaner than the earlier reference sample on fields the institutional-control stack was designed to make explicit. Second, institutional-control semantics appear outside S1/R1/V1, but not with enough coverage to infer full-corpus rates. The selected windows should therefore be read as mature-state deep audits, with P1/B1 as exploratory context.

\subsubsection{4.8 Intervention Timeline}\label{intervention-timeline}

The organization evolved through a sequence of intervention families. The important point is not exact calendar chronology; it is the failure pressure that caused each durable layer to emerge. The organization did not become more controlled through one design insight. Each institutional layer appeared after concrete failure pressure made the previous agent/team abstraction insufficient.

The timeline should be read as institutional growth from failure pressure into durable artifact forms.

{\def\LTcaptype{none} % do not increment counter
\begin{longtable}[]{@{}
  >{\raggedright\arraybackslash}p{(\linewidth - 4\tabcolsep) * \real{0.3333}}
  >{\raggedright\arraybackslash}p{(\linewidth - 4\tabcolsep) * \real{0.3333}}
  >{\raggedright\arraybackslash}p{(\linewidth - 4\tabcolsep) * \real{0.3333}}@{}}
\toprule\noalign{}
\begin{minipage}[b]{\linewidth}\raggedright
Failure pressure
\end{minipage} & \begin{minipage}[b]{\linewidth}\raggedright
Durable artifact form
\end{minipage} & \begin{minipage}[b]{\linewidth}\raggedright
Evidence status
\end{minipage} \\
\midrule\noalign{}
\endhead
\bottomrule\noalign{}
\endlastfoot
Persona drift & role contract and role family & extracted mechanism \\
Chat-as-state & work-record fields & bounded S1 plus mechanism evidence \\
False motion & message freshness and no-reply rule & M1 plus S1 \\
Plausible artifact without authority & verifier decision form & S1 verifier comments \\
Learning trapped in retro & replay and doctrine destination & S1 retro plus mechanism evidence \\
\end{longtable}
}

The full T0-T14 intervention timeline is reported in Appendix I. The main text keeps the compressed failure-to-artifact map because it is the mechanism most relevant to the paper's thesis.

The same history can be summarized as a 0-to-1 evolution arc from weak agent-team forms to institutional-control forms:

{\def\LTcaptype{none} % do not increment counter
\begin{longtable}[]{@{}
  >{\raggedright\arraybackslash}p{(\linewidth - 4\tabcolsep) * \real{0.3333}}
  >{\raggedright\arraybackslash}p{(\linewidth - 4\tabcolsep) * \real{0.3333}}
  >{\raggedright\arraybackslash}p{(\linewidth - 4\tabcolsep) * \real{0.3333}}@{}}
\toprule\noalign{}
\begin{minipage}[b]{\linewidth}\raggedright
Layer
\end{minipage} & \begin{minipage}[b]{\linewidth}\raggedright
Early form
\end{minipage} & \begin{minipage}[b]{\linewidth}\raggedright
Mature institutional form
\end{minipage} \\
\midrule\noalign{}
\endhead
\bottomrule\noalign{}
\endlastfoot
Roles & specialist personas & authority-bounded job contracts and role-family interfaces \\
Workboard & task tracker & institutional state machine \\
Messages & chat or notification & auditable work-transfer trigger \\
Playbooks & optional docs & versioned runtime doctrine \\
Tools/MCP & capability surface & person-role-aware authority surface \\
Harness & task prompt & common employee runtime contract \\
Organization & many agents & routable topology with finality and adaptation \\
\end{longtable}
}

\subsubsection{4.9 Evidence Claim Labels}\label{evidence-claim-labels}

To avoid overclaiming, each result is labeled by evidence type: bounded metric claim, mechanism evidence claim, conceptual claim, design implication, or defined but not estimated in this study. Appendix G lists the allowed wording and evidence basis for major claims.

\subsubsection{4.10 Coding And Evidence Extraction}\label{coding-and-evidence-extraction}

Evidence was extracted from work records, message logs, verifier comments, manager closeouts, replay records, and doctrine changes. Each artifact was coded, when available, for:

\begin{itemize}
\tightlist
\item
  current owner;
\item
  current state;
\item
  exact next action;
\item
  evidence basis;
\item
  role authority;
\item
  verifier decision type;
\item
  no-action state;
\item
  replay state;
\item
  mutation destination;
\item
  human-owner intervention;
\item
  limitation or missing denominator.
\end{itemize}

Claims in the Results section are reported only when supported by bounded denominators or sanitized mechanism packets. Global organization-wide rates require a full corpus extraction and are not claimed in this report.

Each sanitized evidence packet was required to identify the claim, failure mode, intervention, changed durable layer, observed post-change behavior or recurrence check, sanitized artifact, and limitation. This packaging rule was used to keep empirical claims mechanism-grounded rather than asking readers to trust a narrative of improvement.

\subsubsection{4.11 Analysis Strategy}\label{analysis-strategy}

The analysis combines:

\begin{enumerate}
\def\labelenumi{\arabic{enumi}.}
\tightlist
\item
  \textbf{Mechanism tracing:} identifying failure pressure, intervention, changed durable layer, and post-change behavior.
\item
  \textbf{Artifact audit:} checking whether work records, messages, role contracts, playbooks, reviews, closure records, and replay records support the claim.
\item
  \textbf{Case studies:} writing bounded cases where the mechanism is visible.
\item
  \textbf{Bounded metric extraction:} defining organization-level metrics and extracting sample-labeled denominators where the records support them.
\item
  \textbf{Limitations accounting:} recording negative cases, missing recurrence windows, and human-owner boundaries.
\end{enumerate}

\subsection{5. Method: Institutional Control Stack}\label{method-institutional-control-stack}

This section describes the method as a stack of organizational mechanisms. Each mechanism emerged because a simpler agent/team/workflow abstraction failed to control a real operational failure mode.

\subsubsection{5.1 Common Employee Runtime Contract}\label{common-employee-runtime-contract}

The organization used a \textbf{common employee runtime contract} for employee cycles. It is not a prompt contribution in the narrow sense. It is a runtime contract that makes each model invocation answer organizational questions before work can move:

\begin{Shaded}
\begin{Highlighting}[]
\NormalTok{employee cycle =}
\NormalTok{  trigger}
\NormalTok{  + authority}
\NormalTok{  + freshness}
\NormalTok{  + doctrine}
\NormalTok{  + governing record}
\NormalTok{  + evidence}
\NormalTok{  + one authorized action}
\NormalTok{  + audit outcome}
\end{Highlighting}
\end{Shaded}

Freshness means the employee checks whether the trigger is still actionable, already completed, or superseded before taking a step.

The contract frames the model as an employee, not as the system of record. A substantive cycle starts from an inbound message, identifies the governing work record, checks authority, freshness, and evidence, reads relevant doctrine, takes one authorized next step, and leaves a fixed planning/execution stamp with exactly one primary outcome.

A worker cycle was also bounded against backlog sweeping: one trigger had to end in one authorized outcome, not opportunistic cleanup of adjacent work.

The key runtime move is an authority lattice plus evidence demotion. In the studied deployment, authority was ordered from owner instruction and immutable runtime contract, to role text, playbooks, governing work records, and finally messages, comments, retrieved content, and tool outputs as evidence. The employee does not treat all retrieved text as instruction. Role text, playbooks, work records, messages, comments, tool outputs, and copied material answer different institutional questions.

\textbf{Truth-layer separation:} role text owns authority, playbooks own reusable procedure, work records own case truth, messages own activation, and replay owns outcome learning.

{\def\LTcaptype{none} % do not increment counter
\begin{longtable}[]{@{}
  >{\raggedright\arraybackslash}p{(\linewidth - 4\tabcolsep) * \real{0.3333}}
  >{\raggedright\arraybackslash}p{(\linewidth - 4\tabcolsep) * \real{0.3333}}
  >{\raggedright\arraybackslash}p{(\linewidth - 4\tabcolsep) * \real{0.3333}}@{}}
\toprule\noalign{}
\begin{minipage}[b]{\linewidth}\raggedright
Layer
\end{minipage} & \begin{minipage}[b]{\linewidth}\raggedright
Owns
\end{minipage} & \begin{minipage}[b]{\linewidth}\raggedright
Does not own
\end{minipage} \\
\midrule\noalign{}
\endhead
\bottomrule\noalign{}
\endlastfoot
Role text & Employee-specific authority, scope, owned decisions, forbidden decisions, required outputs, and role-specific escalation & Reusable workflow doctrine or current case truth \\
Playbooks & Reusable doctrine, gates, decision criteria, source standards, review standards, and replay standards & Employee authority or ticket-specific state \\
Governing work records & Current case truth: owner, state, evidence, blocker, latest decision, next owner, next action, closure, and replay status & Reusable doctrine or role authority \\
Messages, comments, tool outputs, retrieved content & Evidence, triggers, observations, and handoff requests & Durable institutional authority, role changes, or reusable doctrine \\
\end{longtable}
}

{\def\LTcaptype{none} % do not increment counter
\begin{longtable}[]{@{}
  >{\raggedright\arraybackslash}p{(\linewidth - 4\tabcolsep) * \real{0.3333}}
  >{\raggedright\arraybackslash}p{(\linewidth - 4\tabcolsep) * \real{0.3333}}
  >{\raggedright\arraybackslash}p{(\linewidth - 4\tabcolsep) * \real{0.3333}}@{}}
\toprule\noalign{}
\begin{minipage}[b]{\linewidth}\raggedright
Contract semantic
\end{minipage} & \begin{minipage}[b]{\linewidth}\raggedright
Organizational meaning
\end{minipage} & \begin{minipage}[b]{\linewidth}\raggedright
Failure controlled
\end{minipage} \\
\midrule\noalign{}
\endhead
\bottomrule\noalign{}
\endlastfoot
Authority lattice & Sources are ranked by institutional authority, not recency. & Newer messages, comments, or tool outputs override durable doctrine. \\
Evidence demotion & Retrieved, quoted, or tool-produced text is evidence unless adopted by higher authority. & Instruction-like evidence silently changes role or workflow. \\
Durable instruction layers & Role, playbook, and work record own different kinds of truth. & Case facts leak into doctrine, or doctrine remains trapped in tickets. \\
Governing work record & Every substantive cycle anchors to persistent case truth. & Conversation becomes the system of record. \\
Inbound-message trigger & Work transfer depends on explicit activation. & Assignment fields or mentions silently fail to wake employees. \\
Freshness gate & The trigger is classified as actionable, superseded, or already done. & Duplicate replies, stale writes, and state-sync loops. \\
Planning/execution split & No substantive action before work, authority, state, doctrine, and next step are known. & Execution begins before authority or playbook review. \\
Explicit handoff and silence & Handoffs name owner/action/reply expectation; no-message is valid when no action remains. & Awareness is mistaken for transfer; acknowledgments create false motion. \\
Manager-employee authorized mutation & Durable people and doctrine changes require manager-employee authority inside the human-owner governance boundary. & Tool access becomes institutional authority. \\
Fixed stamps and primary outcomes & The governing record receives planning fields, execution fields, and one outcome class. & Agent traces exist, but organizational state remains unauditable. \\
\end{longtable}
}

The paper treats this contract as a portable organization-protocol concept. It can be implemented with different tools, workboards, and model stacks. Its core contribution is semantic: in this deployment, worker cycles were more auditable and controllable when constrained by authority, state, doctrine, trigger discipline, review, no-action, replay, mutation, and audit rules.

The runtime contract let the organization operate with incomplete doctrine without collapsing into chat. Each cycle still had to locate governing state, respect authority order, demote messages and tool outputs to evidence unless adopted by higher authority, check freshness, take one authorized next step, and leave audit state. As repeated failures exposed missing structure, the manager employee could promote fixes into role contracts, playbooks, work-record fields, message rules, verifier gates, replay doctrine, or tool surfaces inside the human-owner governance boundary. This is evidence of institutional form maturation in one deployment, not proof of unconstrained self-modification or recurrence reduction.

\subsubsection{5.2 Employee Principals}\label{employee-principals}

An AI employee is a persistent principal, not a transient prompt. It has an identity, role, title, status, inbox, tool set, and history. This made it possible to ask: who was asked to act, what role governed the action, what messages did the worker receive, and what authority did the worker hold?

The novelty is not naming agents like employees. The novelty is treating identity as institutional state. An employee can be hired, updated, given or denied tools, messaged, reviewed, and terminated. Those are organizational state transitions, not prompt decorations.

\subsubsection{5.3 Role Contracts}\label{role-contracts}

Roles evolved from humanlike job descriptions into authority contracts. A role contract specifies:

\begin{itemize}
\tightlist
\item
  employee-specific mandate;
\item
  owned decisions;
\item
  forbidden decisions;
\item
  required outputs;
\item
  handoff boundaries;
\item
  escalation rules;
\item
  non-goals.
\end{itemize}

This separated expertise from authority. A research employee could produce an evidence packet without approving execution. A verifier could accept or decline an artifact without placing an order. An executor could act only inside an authorized envelope. A manager could close, reroute, or mutate doctrine, while ordinary employees could only propose durable changes.

A role family also defines an interface: normal upstream artifact, owned transformation, downstream output, and escalation boundary, so hiring or backfilling preserves handoff contracts rather than only adding another capable model.

The general principle is: \textbf{tool capability and subject expertise do not grant institutional authority.}

\subsubsection{5.4 Governing Work Records}\label{governing-work-records}

Every substantive cycle required a governing work record: the persistent artifact that owns case-specific truth. The record holds current owner, state, evidence, blocker, next action, review status, replay linkage, and closure basis.

This addressed a recurring problem: conversation could make work feel complete while the durable work state remained ambiguous. The governing record became the organization's answer to ``what is true now?''

Work records also made parent/child truth explicit. A parent sprint or retro could not be closed merely because some direct children looked complete if a deeper live descendant still held the real blocker or next action. Mature workboard doctrine required recursive descendant truth unless a structural severance was recorded.

\subsubsection{5.5 Workboard As Institutional State Machine}\label{workboard-as-institutional-state-machine}

The workboard became more than a tracker. It was the institutional state machine that turned broad owner intent into executable AI labor. A sprint launch was a compilation step.

Figure 3. Workboard As Institutional State Machine.

\begin{Shaded}
\begin{Highlighting}[]
\NormalTok{owner intent}
\NormalTok{  {-}\textgreater{} live umbrella}
\NormalTok{  {-}\textgreater{} retro umbrella}
\NormalTok{  {-}\textgreater{} child path records}
\NormalTok{  {-}\textgreater{} direct activations}
\NormalTok{  {-}\textgreater{} verifier gates}
\NormalTok{  {-}\textgreater{} close / no{-}action / replay}
\end{Highlighting}
\end{Shaded}

Takeaway: broad owner intent becomes routable, verifiable, closeable work only after it is compiled into durable workboard state.

The live umbrella owned the active frontier. The retro umbrella owned later outcome truth and learning obligations. Child records owned episode truth: current owner, gate state, evidence, blocker, next action, and closure or no-action basis. Direct messages activated one employee for one child path rather than turning a broad sprint request into a general inbox sweep.

This made the workboard a frontier controller. A parent record could not truthfully close while a child still governed live work. A parent also should not remain active after the child frontier was truthfully complete. Replay obligations moved to the retro surface so live records did not stay fake-active while waiting for delayed outcome truth.

The portable design claim is that high-skill AI labor benefits from a state substrate that compiles vague intent into small, routable, verifiable work objects with explicit activation and closure semantics.

\subsubsection{5.6 Work Messages}\label{work-messages}

Messages became work triggers rather than chat. A valid work message names the work item, exact ask, exact next action, and whether reply is needed.

This design emerged because tracker assignment, comments, and mentions were insufficient to wake AI employees reliably. It later evolved further because messages themselves can become stale or duplicate. The mature protocol added inbox pre-checks, \texttt{ACTIONABLE}, \texttt{SUPERSEDED}, and \texttt{ALREADY\_DONE} classifications, no-reply suppression, and limits on self-messages.

Inbox transparency was scoped. Employees checked newer relevant messages for the same work item to classify the trigger as actionable, superseded, or already done. They were not meant to use inbox access as a general monitoring surface, second workboard, or backlog sweep.

One counterintuitive result is that correct silence became a control behavior. If the governing work record is truthful and no recipient must act, sending a completion note can create false motion. Message evidence shows repeated manager closeouts where \texttt{Reply\ needed:\ NO} was honored and no direct message was sent.

The system also avoided awareness-only copies because extra messages are not harmless notifications. For an AI employee, every inbound message can become a work trigger unless the institution gives it explicit no-action or no-reply semantics.

\subsubsection{5.7 Playbooks As Operating Doctrine}\label{playbooks-as-operating-doctrine}

Playbooks are reusable doctrine, not optional retrieved context. They own workflow logic, domain-specific acceptance criteria, required inputs and outputs, source truth, replay standards, and recurring review steps.

Playbooks were not simply pre-authored instructions. They became the durable landing zone for reusable lessons discovered during work. Case-specific facts remained in work records; reusable procedure moved into playbooks; authority remained in role text. This placement discipline is what turned ad hoc correction into institutional learning.

This differs from generic documentation in three ways:

\begin{enumerate}
\def\labelenumi{\arabic{enumi}.}
\tightlist
\item
  Runtime doctrine must be listed and materially relevant playbooks must be read before substantive action.
\item
  Reusable learning must be promoted to the correct durable layer rather than left in a ticket or message thread.
\item
  Playbook changes are institutional mutations, not ad hoc prompt edits.
\end{enumerate}

Because role text sits above playbooks in the authority stack, reusable procedures should not be copied wholesale into roles. Stale workflow embedded in role text can override newer playbook doctrine. Mature role design keeps roles focused on authority, interface, outputs, and escalation, while playbooks own detailed operating procedure.

The doctrine became more procedural exactly where failures recurred: source freshness, packet completeness, activation, no-reply handoffs, replay linkage, and doctrine editing.

The organization also learned that long playbooks and role texts are necessary for explicit model control but risky to rewrite wholesale. This motivated targeted old-string/new-string doctrine edits, which preserve local explicitness and reduce accidental compression or restyling. This report treats targeted doctrine patching as mechanism evidence; quantitative edit-error, context-cost, and recurrence values are not reported as bounded values.

\subsubsection{5.8 Evidence Surfaces And Decision Packets}\label{evidence-surfaces-and-decision-packets}

High-skill work cannot be reviewed if it is only plausible prose. Weather and market research therefore had to become verifier-reconstructible packets: source basis, settlement mapping, uncertainty, scenario reasoning, probability or decision structure, confidence, contradiction checks, freshness, and no-action triggers.

The generalizable mechanism is not weather-specific. In technical AI organizations, downstream review benefits from artifacts a second employee can verify without redoing the producer's entire job. The producer's artifact should expose enough evidence and assumptions for a verifier to accept, reject, repair, reroute, or decline.

\subsubsection{5.9 Verifier Topology}\label{verifier-topology}

The organization used verifier topology rather than verifier prompting. The basic topology was:

\begin{enumerate}
\def\labelenumi{\arabic{enumi}.}
\tightlist
\item
  Creator produces an artifact.
\item
  Functional verifier accepts, declines, repairs, or reroutes.
\item
  Executor acts only inside accepted authority.
\item
  Manager verifies closure and durable learning.
\end{enumerate}

Figure 4. Creator -\textgreater{} Verifier -\textgreater{} Executor -\textgreater{} Manager -\textgreater{} Replay.

\begin{Shaded}
\begin{Highlighting}[]
\NormalTok{creator artifact}
\NormalTok{  {-}\textgreater{} functional verifier decision}
\NormalTok{  {-}\textgreater{} authorized execution or no{-}action}
\NormalTok{  {-}\textgreater{} manager closure decision}
\NormalTok{  {-}\textgreater{} replay against outcome truth}
\end{Highlighting}
\end{Shaded}

Takeaway: review becomes an authority-bearing route through state decisions, not a prose critique step.

This made verification an authority-bearing state transition. A verifier message did not ask for ``thoughts.'' It asked for a governing decision. Message evidence showed verifier responses ending in approval, decline, repair handoff, reroute, or manager review. Some approvals changed the downstream executable envelope rather than rubber-stamping the upstream proposal. Some declines explicitly withheld execution authorization even when the packet quality passed.

Verifier topology separates artifact review from execution authorization. A packet can be reconstructible yet still fail the action gate. In those cases, the verifier can decline execution, request repair, reroute the work, request manager review, or narrow the executable envelope.

A sanitized verifier decision can be represented without exposing market strategy or raw operational records:

\begin{Shaded}
\begin{Highlighting}[]
\NormalTok{Work item: [redacted market{-}path episode]}
\NormalTok{Decision: Repair required}
\NormalTok{Reason: settlement mapping incomplete; source freshness not established}
\NormalTok{Next owner: research employee}
\NormalTok{Next action: refresh source basis and resubmit packet}
\NormalTok{Execution authorized: no}
\end{Highlighting}
\end{Shaded}

\subsubsection{5.10 No-Action States}\label{no-action-states}

No-action became first-class. Valid outcomes included no-send, decline, repair, stand-down, blocked, already-done, explicit zero, no replay needed, and case-only learning. These outcomes are not failures by default. They are often the organization's correct output.

Explicit-zero is a verified no-action state: the organization records that no live child path, downstream owner, or executable action remains. It is completed work only because the governing record is truthful and accepted, not because an outcome was proven correct.

Figure 5. Work Item Lifecycle.

\begin{Shaded}
\begin{Highlighting}[]
\NormalTok{trigger received}
\NormalTok{  {-}\textgreater{} freshness state: actionable / superseded / already done}
\NormalTok{  {-}\textgreater{} governing record identified}
\NormalTok{  {-}\textgreater{} planning gate: authority + doctrine + evidence}
\NormalTok{  {-}\textgreater{} primary outcome:}
\NormalTok{       completed / handoff required / blocked / needs review / needs clarification}
\NormalTok{  {-}\textgreater{} closure / no{-}action / replay / mutation review}
\end{Highlighting}
\end{Shaded}

Takeaway: the lifecycle treats stopping, handoff, review, replay, and mutation review as explicit institutional states.

This matters because AI systems tend to reward visible action. In high-skill work, value often comes from refusing unsupported action. The organization therefore had to make stopping auditable: why no action was correct, who had authority to decide, whether review was required, whether replay was needed, and who owned the next step if any.

Each substantive handling cycle ended in exactly one primary outcome:

{\def\LTcaptype{none} % do not increment counter
\begin{longtable}[]{@{}
  >{\raggedright\arraybackslash}p{(\linewidth - 2\tabcolsep) * \real{0.5000}}
  >{\raggedright\arraybackslash}p{(\linewidth - 2\tabcolsep) * \real{0.5000}}@{}}
\toprule\noalign{}
\begin{minipage}[b]{\linewidth}\raggedright
Primary outcome
\end{minipage} & \begin{minipage}[b]{\linewidth}\raggedright
Meaning
\end{minipage} \\
\midrule\noalign{}
\endhead
\bottomrule\noalign{}
\endlastfoot
Completed & Work is done, including valid no-change or no-action completion \\
Handoff required & Another owner must act \\
Blocked & A real dependency prevents progress \\
Needs review & A verifier or manager must decide \\
Needs clarification & A narrow missing fact or ambiguity blocks correct action \\
\end{longtable}
}

\subsubsection{5.11 Replay And Institutional Mutation}\label{replay-and-institutional-mutation}

Replay connects outcome truth to organizational learning. A replay record reconstructs the original path, active role/playbook context, decision artifact, action or no-action, later outcome truth, process correctness, outcome correctness, and learning destination.

The important replay detail is historical reconstruction. Replay asks what role contract and playbook versions governed the employee at the time, not what the latest doctrine says now. This prevents the organization from judging old work against rules that were not in force at the time and makes expected-versus-realized review possible.

Learning destinations include:

\begin{itemize}
\tightlist
\item
  role text;
\item
  playbook;
\item
  tool access;
\item
  staffing;
\item
  topology;
\item
  workboard rule;
\item
  message rule;
\item
  case-only record.
\end{itemize}

This is the adaptive institution loop:

\begin{quote}
evidence -\textgreater{} classification -\textgreater{} authorized mutation -\textgreater{} activation -\textgreater{} recurrence check.
\end{quote}

Reflection is insufficient because it can stay private to one model invocation. Institutional learning is complete only when the reusable fix lands in the correct durable layer or is intentionally classified as case-only.

A sanitized replay record can be represented as:

\begin{Shaded}
\begin{Highlighting}[]
\NormalTok{Original work path: [redacted market{-}path episode]}
\NormalTok{Original decision artifact: evidence packet plus verifier decision}
\NormalTok{Role version at decision time: [redacted role contract version]}
\NormalTok{Playbook version at decision time: [redacted review/playbook version]}
\NormalTok{Action state: action / no{-}action / stand{-}down}
\NormalTok{Outcome truth: later observed result or settlement state}
\NormalTok{Expected{-}vs{-}realized review: process{-}correct / process{-}gap / outcome{-}gap / unclear}
\NormalTok{Learning destination: playbook / role / tool / workboard / message / topology / case{-}only}
\NormalTok{Mutation status: proposed / applied / rejected / deferred / case{-}only}
\end{Highlighting}
\end{Shaded}

\subsubsection{5.12 Person-Role-Aware Tool Surface}\label{person-role-aware-tool-surface}

The tool surface was split into two planes:

\begin{enumerate}
\def\labelenumi{\arabic{enumi}.}
\tightlist
\item
  \textbf{Work-tool plane:} domain tools used to gather evidence, inspect markets, compute, or execute.
\item
  \textbf{Institutional-control plane:} tools for employees, messages, profiles, playbooks, role edits, tool edits, hiring, termination, and doctrine mutation.
\end{enumerate}

The important idea is not the particular implementation. It is that organization-grade tool use must be interpreted through employee identity, role authority, governing work state, and doctrine. MCP-style tools expose capability. An organization layer decides when that capability may be used and what durable state change it creates.

In this sense, the tool surface functioned as an executable authority model. An authenticated employee principal entered the organization through a caller-dependent surface: ordinary employees could read profiles, inspect inboxes, read playbooks, message coworkers, and perform assigned work; the manager-employee surface exposed scoped institutional mutation capabilities such as role edits, tool-set changes, hiring, termination, and playbook mutation. The human owner remained the outer governance authority. Profile, role, tool-set, inbox, work-record, and playbook reads were intentionally separated because each surface answers a different authority question.

The table abstracts the public-facing surface into institutional semantics; it is not a complete tool inventory.

{\def\LTcaptype{none} % do not increment counter
\begin{longtable}[]{@{}
  >{\raggedright\arraybackslash}p{(\linewidth - 6\tabcolsep) * \real{0.2500}}
  >{\raggedright\arraybackslash}p{(\linewidth - 6\tabcolsep) * \real{0.2500}}
  >{\raggedright\arraybackslash}p{(\linewidth - 6\tabcolsep) * \real{0.2500}}
  >{\raggedright\arraybackslash}p{(\linewidth - 6\tabcolsep) * \real{0.2500}}@{}}
\toprule\noalign{}
\begin{minipage}[b]{\linewidth}\raggedright
Surface
\end{minipage} & \begin{minipage}[b]{\linewidth}\raggedright
Regular employee
\end{minipage} & \begin{minipage}[b]{\linewidth}\raggedright
Manager employee / human owner
\end{minipage} & \begin{minipage}[b]{\linewidth}\raggedright
Institutional meaning
\end{minipage} \\
\midrule\noalign{}
\endhead
\bottomrule\noalign{}
\endlastfoot
Profile / role reads & yes & yes & identity and authority are inspectable \\
Inbox / message tools & yes & yes & communication is work activation \\
Playbook reads & yes & yes & doctrine is readable at runtime \\
Role/tool/playbook mutation & no & yes & durable institutional change is privileged \\
Hiring / termination & no & yes & employees are institutional principals \\
\end{longtable}
}

Tool outputs were also shaped for institutional use: list and inbox surfaces were paginated, outputs were minimized, message status was exposed, and role/playbook version reads supported audit and replay.

For communication, the institutional tool surface exposed sender identity, message status, timestamps, and response/failure fields so inbox reads could support obligation freshness rather than become a general chat archive.

A sanitized authority example is:

\begin{Shaded}
\begin{Highlighting}[]
\NormalTok{Employee state: can inspect evidence and draft a recommendation}
\NormalTok{Work{-}tool assignment: provides the tools needed for this role\textquotesingle{}s assigned work}
\NormalTok{Role authority: cannot authorize downstream execution or closure}
\NormalTok{Required gate: verifier acceptance and authorized execution envelope}
\NormalTok{Valid outcome: record needs{-}review or handoff state; message the authorized next owner}
\NormalTok{Invalid outcome: treat evidence access or work{-}tool access as permission to execute}
\end{Highlighting}
\end{Shaded}

This is the concrete version of ``tool access is not authority.'' The institutional surface assigns the tools needed for a role's work, while the valid state mutation still depends on the employee principal, role contract, governing work record, verifier gate, and manager-employee authority boundary.

\subsection{6. Evaluation Framework}\label{evaluation-framework}

The evaluation asks:

\begin{quote}
Did the organization produce accountable, authorized, evidence-backed, outcome-reconciled work, and could reliability-relevant failures be audited, bounded, or repaired?
\end{quote}

It deliberately avoids using profit, number of employees, number of tools, prompt length, or message volume as primary evidence. Those are context. They do not prove institutional control.

\subsubsection{6.1 Metric Families}\label{metric-families}

The metrics are not task-success metrics. They are institution-state metrics: they ask whether the organization made the right work truth visible and actionable.

Metric status is reported in four tiers: \texttt{defined} means the metric is conceptually specified; \texttt{extractable} means source artifacts are known but no denominator-labeled value is reported; \texttt{mechanism\ evidence\ extracted} means sanitized cases support the mechanism; \texttt{bounded\ value\ extracted} means Section 7 reports a numerator, denominator, sample label, and window. No denominator means no bounded value.

{\def\LTcaptype{none} % do not increment counter
\begin{longtable}[]{@{}
  >{\raggedright\arraybackslash}p{(\linewidth - 6\tabcolsep) * \real{0.2500}}
  >{\raggedright\arraybackslash}p{(\linewidth - 6\tabcolsep) * \real{0.2500}}
  >{\raggedright\arraybackslash}p{(\linewidth - 6\tabcolsep) * \real{0.2500}}
  >{\raggedright\arraybackslash}p{(\linewidth - 6\tabcolsep) * \real{0.2500}}@{}}
\toprule\noalign{}
\begin{minipage}[b]{\linewidth}\raggedright
Metric
\end{minipage} & \begin{minipage}[b]{\linewidth}\raggedright
What it checks
\end{minipage} & \begin{minipage}[b]{\linewidth}\raggedright
Evidence/status
\end{minipage} & \begin{minipage}[b]{\linewidth}\raggedright
Main limitation
\end{minipage} \\
\midrule\noalign{}
\endhead
\bottomrule\noalign{}
\endlastfoot
Ownership validity & owner and next action/blocker & P1/B1 spot-checks; S1 closeouts; bounded value extracted & P1/B1 shallow; S1 selected \\
Authority compliance & role, doctrine, owner/manager, and state permission & mechanism evidence extracted & mechanism evidence only \\
Evidence sufficiency & packet reviewability & mechanism evidence extracted & no denominator \\
Verifier gate quality & stateful accept, repair, reroute, decline, or review & mechanism evidence extracted & lower bounds, not rates \\
Closure validity & authority-accepted final state & S1 closeouts; bounded value extracted & selected window \\
Workboard frontier validity & parent, child, activation, closure, and replay state agreement & mechanism evidence extracted & duplicate-parent case is not denominator-coded \\
No-action quality & evidence basis and no-replay treatment & S1 no-execution paths; bounded value extracted & selected window; not correctness \\
Message trigger auditability & exact ask, reply flag, and stale suppression & M1 packets, S1 closeouts, and S1 live umbrella; mechanism evidence extracted & S1 activation coverage comes from workboard/umbrella evidence, not a full raw inbox census \\
Replay completion & replay or no-replay rationale & S1 replay/final records; bounded value extracted & no recurrence rate \\
Topology routability & lane, chain, tool surface, and activation path identify owner & mechanism evidence extracted & mechanism evidence only \\
Role-family activation readiness & required role, work-tool, routing, and activation agreement & mechanism evidence extracted & mechanism-level evidence only; no bounded estimate reported \\
Durable mutation latency & time from learning to durable-layer change & defined & defined but not estimated in this study \\
Recurrence after mutation & same failure after activated change & defined & defined but not estimated in this study \\
Human intervention burden & owner or human repair frequency & mechanism evidence extracted & not global burden \\
\end{longtable}
}

Appendix H records detailed evidence status and coding guidance for metrics this study defines but does not estimate. The main Results section keeps only denominator-labeled counts and short evidence-basis notes.

\subsection{7. Results}\label{results}

This section reports defensible findings from extracted evidence. It does not invent aggregate organization-wide rates. Quantitative rows use bounded samples: the primary window is \texttt{S1}, a 20-child live sprint paired with its retro; \texttt{R1} and \texttt{V1} are used only as labeled supporting windows.

The Results section measures institutional auditability and controllability, not trading performance, forecasting accuracy, or global decision quality. Counts are reported only for bounded windows with denominator-labeled extraction. Where evidence is qualitative or mechanism-level, the text does not infer aggregate rates.

Results at a glance:

{\def\LTcaptype{none} % do not increment counter
\begin{longtable}[]{@{}
  >{\raggedright\arraybackslash}p{(\linewidth - 6\tabcolsep) * \real{0.2500}}
  >{\raggedright\arraybackslash}p{(\linewidth - 6\tabcolsep) * \real{0.2500}}
  >{\raggedright\arraybackslash}p{(\linewidth - 6\tabcolsep) * \real{0.2500}}
  >{\raggedright\arraybackslash}p{(\linewidth - 6\tabcolsep) * \real{0.2500}}@{}}
\toprule\noalign{}
\begin{minipage}[b]{\linewidth}\raggedright
Result family
\end{minipage} & \begin{minipage}[b]{\linewidth}\raggedright
Main bounded evidence
\end{minipage} & \begin{minipage}[b]{\linewidth}\raggedright
What it shows
\end{minipage} & \begin{minipage}[b]{\linewidth}\raggedright
What it does not show
\end{minipage} \\
\midrule\noalign{}
\endhead
\bottomrule\noalign{}
\endlastfoot
Ownership and next action & P1: 21/30 owner, 18/30 next; S1: 20/20 both; B1: 62/75 owner, 51/75 next & mature S1 made ownership and next-step state explicit & full-corpus rate or correctness \\
Closure & S1: 20/20 manager-reviewed; 0/20 complete-but-open after final closeout & completion was separated from accepted closure & global closure-validity rate \\
No-action & S1: 6/6 no-execution paths with explicit rationale/no-replay treatment & stopping became an auditable output & outcome correctness \\
Replay & S1: 14/14 executed paths replay-closed & delayed outcome truth stayed linked to live work & recurrence reduction \\
Communication & S1: direct activation recorded for 20/20 child records; M1: stale/no-change suppression examples & handoff and silence became auditable & inbox-wide loop rate \\
\end{longtable}
}

Across extracted mechanism packets, institutional state made downstream decisions and transitions auditable in ways that ordinary record-completeness tables do not fully show:

{\def\LTcaptype{none} % do not increment counter
\begin{longtable}[]{@{}
  >{\raggedright\arraybackslash}p{(\linewidth - 10\tabcolsep) * \real{0.1667}}
  >{\raggedright\arraybackslash}p{(\linewidth - 10\tabcolsep) * \real{0.1667}}
  >{\raggedright\arraybackslash}p{(\linewidth - 10\tabcolsep) * \real{0.1667}}
  >{\raggedright\arraybackslash}p{(\linewidth - 10\tabcolsep) * \real{0.1667}}
  >{\raggedright\arraybackslash}p{(\linewidth - 10\tabcolsep) * \real{0.1667}}
  >{\raggedright\arraybackslash}p{(\linewidth - 10\tabcolsep) * \real{0.1667}}@{}}
\toprule\noalign{}
\begin{minipage}[b]{\linewidth}\raggedright
Pillar
\end{minipage} & \begin{minipage}[b]{\linewidth}\raggedright
Mechanism
\end{minipage} & \begin{minipage}[b]{\linewidth}\raggedright
Auditable decision or transition
\end{minipage} & \begin{minipage}[b]{\linewidth}\raggedright
Evidence type
\end{minipage} & \begin{minipage}[b]{\linewidth}\raggedright
Claim type
\end{minipage} & \begin{minipage}[b]{\linewidth}\raggedright
Limitation
\end{minipage} \\
\midrule\noalign{}
\endhead
\bottomrule\noalign{}
\endlastfoot
State & Manager closure & Apparently complete work was kept open, rerouted, or closed only after final disposition. & S1 closeout audit & bounded metric & selected mature window \\
Authority & Verifier gate & Approval, repair, decline, reroute, and changed execution envelope became state decisions. & S1 verifier comments & mechanism evidence & lower bounds, not rates \\
Communication & Message freshness & No-change and superseded triggers were suppressed before action. & M1 packet examples & extracted mechanism & not inbox census \\
State & Workboard compilation & Broad intent became child paths, gates, activations, and replay obligations. & S1/R1/V1 sprint structure & mechanism evidence & no corpus-wide rate \\
Learning & Replay classification & Live closure, no-replay rationale, and doctrine follow-up were separated. & S1 replay and retro records & bounded metric & no recurrence rate \\
Authority & Activation-gated role-family expansion & Activation waited until role, tool, routing, and activation surfaces aligned. & Case I mechanism evidence & extracted mechanism & no recurrence or safety claim \\
\end{longtable}
}

\subsubsection{7.1 Finding 1: State Control Made Work Truth Auditable}\label{finding-1-state-control-made-work-truth-auditable}

\textbf{Evidence basis:} bounded metric audit.

State control converted ambiguous conversational residue into auditable institutional state. Work that looked substantively complete still required an accepted final state, named dependency, reroute, or explicit no-action treatment on the governing record.

A bounded runtime-contract compliance audit on the S1 final manager-closeout sample found that 20/20 final closeout records included fixed planning and execution stamps, freshness state, materially relevant playbooks read, exactly one primary outcome, next owner/action fields, and a messages-sent field. In the same sample, 20/20 final manager closeouts sent no direct message when no downstream action remained, and 6/6 final no-execution paths were audited as valid no-action with explicit no-replay treatment. These are not global rates; they show that the mature contract made closure, no-action, and silence visible on the governing record.

These presence metrics measure record completeness and auditability, not decision correctness. A field was counted when it was present in the fixed record structure; the count does not by itself prove the underlying decision was correct.

{\def\LTcaptype{none} % do not increment counter
\begin{longtable}[]{@{}
  >{\raggedright\arraybackslash}p{(\linewidth - 4\tabcolsep) * \real{0.3333}}
  >{\raggedright\arraybackslash}p{(\linewidth - 4\tabcolsep) * \real{0.3333}}
  >{\raggedright\arraybackslash}p{(\linewidth - 4\tabcolsep) * \real{0.3333}}@{}}
\toprule\noalign{}
\begin{minipage}[b]{\linewidth}\raggedright
Runtime-contract compliance field
\end{minipage} & \begin{minipage}[b]{\linewidth}\raggedright
S1 final closeout denominator
\end{minipage} & \begin{minipage}[b]{\linewidth}\raggedright
Count
\end{minipage} \\
\midrule\noalign{}
\endhead
\bottomrule\noalign{}
\endlastfoot
Planning and execution stamps present & 20 final closeout records & 20/20 \\
Freshness state, playbooks read, primary outcome, next owner/action, and messages-sent fields present & 20 final closeout records & 20/20 \\
Final closeouts with no direct message sent when no action remained & 20 final closeout records & 20/20 \\
No-execution paths audited as valid no-action & 6 final no-execution paths & 6/6 \\
\end{longtable}
}

The same evidence supports an institutional distinction between completion and closure:

\begin{itemize}
\tightlist
\item
  \textbf{completion}: a worker's assigned output is finished;
\item
  \textbf{closure}: an authorized verifier has accepted the state and recorded final truth.
\end{itemize}

Work records showed that completion and closure diverged. Some work appeared substantively done but remained open, review-ready, or structurally incomplete. The manager-verification intervention required manager-touched records to end as one of: closed, kept open for a named dependency, rejected with the missing piece, or rerouted to the next owner.

{\def\LTcaptype{none} % do not increment counter
\begin{longtable}[]{@{}
  >{\raggedright\arraybackslash}p{(\linewidth - 4\tabcolsep) * \real{0.3333}}
  >{\raggedright\arraybackslash}p{(\linewidth - 4\tabcolsep) * \real{0.3333}}
  >{\raggedright\arraybackslash}p{(\linewidth - 4\tabcolsep) * \real{0.3333}}@{}}
\toprule\noalign{}
\begin{minipage}[b]{\linewidth}\raggedright
Metric
\end{minipage} & \begin{minipage}[b]{\linewidth}\raggedright
Status
\end{minipage} & \begin{minipage}[b]{\linewidth}\raggedright
Value
\end{minipage} \\
\midrule\noalign{}
\endhead
\bottomrule\noalign{}
\endlastfoot
Complete-but-open exceptions after final manager closeout & S1 live child tickets & 0/20 observed \\
Manager touches with final disposition after intervention & S1 live child tickets & 20/20 final live dispositions manager-reviewed; 2 pre-final manager reroutes bounded premature finality \\
Review-ready records left without state transition & S1 live child tickets & 0/20 observed after final manager verification \\
\end{longtable}
}

\textbf{Limitation.} These are record-completeness and auditability metrics from a selected mature window, not correctness or full deployment rates.

\subsubsection{7.2 Finding 2: Authority Control Turned Review Into State Decisions}\label{finding-2-authority-control-turned-review-into-state-decisions}

\textbf{Evidence basis:} mechanism evidence with bounded lower-bound counts.

Verifier evidence shows that review messages did not ask for generic critique. They asked the verifier to make a governing decision. The verifier could approve, decline, return repair items, reroute, or request manager review. Approval could also modify the executable envelope rather than rubber-stamping the upstream proposal. This made authority a topology property rather than a tone of reviewer feedback.

Verifier topology separated two questions that agent systems often collapse: ``Is the artifact reconstructible?'' and ``Is downstream action authorized?'' A packet could be reviewable while still failing the action gate, causing a verifier to decline execution, request repair, reroute, request manager review, or narrow the executable envelope.

Rows using ``at least'' are conservative lower bounds from extracted comments, not event-rate estimates over the full deployment.

{\def\LTcaptype{none} % do not increment counter
\begin{longtable}[]{@{}
  >{\raggedright\arraybackslash}p{(\linewidth - 4\tabcolsep) * \real{0.3333}}
  >{\raggedright\arraybackslash}p{(\linewidth - 4\tabcolsep) * \real{0.3333}}
  >{\raggedright\arraybackslash}p{(\linewidth - 4\tabcolsep) * \real{0.3333}}@{}}
\toprule\noalign{}
\begin{minipage}[b]{\linewidth}\raggedright
Verifier outcome
\end{minipage} & \begin{minipage}[b]{\linewidth}\raggedright
Status
\end{minipage} & \begin{minipage}[b]{\linewidth}\raggedright
Observed count or disposition
\end{minipage} \\
\midrule\noalign{}
\endhead
\bottomrule\noalign{}
\endlastfoot
Approve & S1 functional-verifier comments & At least 16 approval events observed \\
Repair required & S1 functional-verifier comments & At least 10 repair-required decisions across at least 9 child paths \\
Decline/no execution & S1 final live child dispositions & 6/20 final no-execution paths \\
Reroute & S1 comments & 4 observed reroutes: 2 execution-stage and 2 manager-stage \\
Needs manager review & S1 live child tickets & 20/20 final dispositions routed through manager verification \\
Approval changed upstream proposed action & S1 comments & 3 clear changed-envelope approvals \\
\end{longtable}
}

Authority control also constrained execution. A proposed action was not executable merely because a worker generated it or a tool could perform it. The action had to pass through role authority, verifier gates, and manager closure. This is the paper's main distinction from agent-level guardrail framing: the gate is attached to an organizational state transition, not only to a model output or tool call.

\textbf{Limitation.} The verifier counts are conservative observed lower bounds in extracted comments, not full verifier-event rates.

\subsubsection{7.3 Finding 3: False-Motion Control Made Stale, Duplicate, And Non-Material Work Suppressible}\label{finding-3-false-motion-control-made-stale-duplicate-and-non-material-work-suppressible}

\textbf{Evidence basis:} mechanism evidence plus bounded closeout counts.

Communication control made false motion suppressible. Duplicate state sync, no-change updates, waiting loops, fake-active live work, stale parent umbrellas, and broad sprint triggers can all look like activity while failing to change organizational truth; the message protocol and freshness checks, using scoped inbox context, made no-reply, stale suppression, and narrowed activation auditable outcomes.

{\def\LTcaptype{none} % do not increment counter
\begin{longtable}[]{@{}
  >{\raggedright\arraybackslash}p{(\linewidth - 4\tabcolsep) * \real{0.3333}}
  >{\raggedright\arraybackslash}p{(\linewidth - 4\tabcolsep) * \real{0.3333}}
  >{\raggedright\arraybackslash}p{(\linewidth - 4\tabcolsep) * \real{0.3333}}@{}}
\toprule\noalign{}
\begin{minipage}[b]{\linewidth}\raggedright
Metric
\end{minipage} & \begin{minipage}[b]{\linewidth}\raggedright
Status
\end{minipage} & \begin{minipage}[b]{\linewidth}\raggedright
Value
\end{minipage} \\
\midrule\noalign{}
\endhead
\bottomrule\noalign{}
\endlastfoot
Manager-review messages with no-reply expectation & S1 closeout records plus message packet sample & Observed as a closeout pattern; no global inbox rate claimed \\
Closeouts with no direct message sent & S1 live child manager closeouts & 20/20 \\
Direct activation coverage & S1 live umbrella & 20/20 child records direct-activated \\
Self-trigger no-change suppressions & 260-row inbox packet sample & 1 sanitized packet \\
Superseded/stale unauthorized-message suppressions & 260-row inbox packet sample & 2 sanitized packets \\
Broad sprint triggers narrowed into child activations & S1/R1/V1 sprint structure & Mechanism observed; no corpus-wide rate claimed \\
\end{longtable}
}

\textbf{Limitation.} These examples show false-motion suppression in bounded samples; they do not establish loop elimination.

The S1 live umbrella recorded direct activation for 20/20 child records; this supports the claim that handoff was treated as an activation event, not merely assignment metadata. The activation coverage comes from workboard/umbrella evidence, not a full raw inbox census.

\subsubsection{7.4 Finding 4: Learning Control Separated Stopping, Replay, And Doctrine Change}\label{finding-4-learning-control-separated-stopping-replay-and-doctrine-change}

\textbf{Evidence basis:} bounded metric audit.

Learning control separated stopping, live closure, delayed outcome truth, and durable doctrine change. In the S1 no-execution paths, no-action was not treated as missing work: it was made reviewable through evidence basis, authority basis, and explicit replay or no-replay treatment. In the strongest form, explicit zero is verified non-action: the organization records that no live path remains and manager closure accepts that state.

{\def\LTcaptype{none} % do not increment counter
\begin{longtable}[]{@{}
  >{\raggedright\arraybackslash}p{(\linewidth - 4\tabcolsep) * \real{0.3333}}
  >{\raggedright\arraybackslash}p{(\linewidth - 4\tabcolsep) * \real{0.3333}}
  >{\raggedright\arraybackslash}p{(\linewidth - 4\tabcolsep) * \real{0.3333}}@{}}
\toprule\noalign{}
\begin{minipage}[b]{\linewidth}\raggedright
No-action metric
\end{minipage} & \begin{minipage}[b]{\linewidth}\raggedright
Status
\end{minipage} & \begin{minipage}[b]{\linewidth}\raggedright
Value
\end{minipage} \\
\midrule\noalign{}
\endhead
\bottomrule\noalign{}
\endlastfoot
No-action records by type & S1 live child tickets & 6/20 final no-execution paths: 4 verifier declines and 2 no-send/weak-edge decisions \\
No-action records with evidence basis & S1 no-execution paths & 6/6 \\
No-action outcomes later revisited or classified & S1 comments and retro & 2 no-action/decline states were rerouted before final closure; final no-execution paths were classified no-replay-required \\
\end{longtable}
}

Replay separated live-work closure from delayed outcome truth. An executed or no-action path could be closed operationally while a replay child or retro surface remained open for later settlement, fill, cancellation, or outcome review.

The important learning mechanism was destination classification. A replay could conclude that learning belonged in a playbook, role, tool surface, staffing/topology, workboard rule, message protocol, or only the case record. This controlled two common failure modes: over-promoting one-off facts into doctrine and under-promoting repeated failures into dead comments.

{\def\LTcaptype{none} % do not increment counter
\begin{longtable}[]{@{}
  >{\raggedright\arraybackslash}p{(\linewidth - 4\tabcolsep) * \real{0.3333}}
  >{\raggedright\arraybackslash}p{(\linewidth - 4\tabcolsep) * \real{0.3333}}
  >{\raggedright\arraybackslash}p{(\linewidth - 4\tabcolsep) * \real{0.3333}}@{}}
\toprule\noalign{}
\begin{minipage}[b]{\linewidth}\raggedright
Replay metric
\end{minipage} & \begin{minipage}[b]{\linewidth}\raggedright
Status
\end{minipage} & \begin{minipage}[b]{\linewidth}\raggedright
Value
\end{minipage} \\
\midrule\noalign{}
\endhead
\bottomrule\noalign{}
\endlastfoot
Ordered paths with replay complete & S1 replay tickets & 14/14 closed under the retro umbrella \\
No-execution paths with explicit no-replay rationale & S1 no-execution paths & 6/6 \\
Replay findings promoted to durable mutation & S1 retro closeout & 0 immediate durable patches; 1 research/doctrine follow-up artifact \\
Replay findings classified case-only & S1 retro closeout & Ticket-level replay records were closed, but per-ticket destination split was not separately coded \\
\end{longtable}
}

\textbf{Limitation.} The no-action and replay counts show reviewability and closure treatment, not outcome correctness or recurrence reduction.

\subsubsection{7.5 Finding 5: Topology Control Made The Organization Routable}\label{finding-5-topology-control-made-the-organization-routable}

\textbf{Evidence basis:} mechanism evidence.

The organization did not gain auditable control merely by adding more agents. It gained routable institutional control when agent seats were placed into topology: role families, pods or lanes, default downstream chains, verifier and executor separation, manager finality, capacity/overflow concepts, production/research separation, and synchronized truth surfaces.

This finding is reported as mechanism evidence, not a topology-performance result. The observed pattern is that structural change was only operational when role text, profile or roster truth, tool access, work records, and activation messages agreed. A role could exist in design but remain ineffective if it lacked the right tool surface or activation path. A tool-capable employee could still lack authority if the role, verifier gate, or governing record did not authorize action.

{\def\LTcaptype{none} % do not increment counter
\begin{longtable}[]{@{}
  >{\raggedright\arraybackslash}p{(\linewidth - 6\tabcolsep) * \real{0.2500}}
  >{\raggedright\arraybackslash}p{(\linewidth - 6\tabcolsep) * \real{0.2500}}
  >{\raggedright\arraybackslash}p{(\linewidth - 6\tabcolsep) * \real{0.2500}}
  >{\raggedright\arraybackslash}p{(\linewidth - 6\tabcolsep) * \real{0.2500}}@{}}
\toprule\noalign{}
\begin{minipage}[b]{\linewidth}\raggedright
Topology mechanism
\end{minipage} & \begin{minipage}[b]{\linewidth}\raggedright
Control function
\end{minipage} & \begin{minipage}[b]{\linewidth}\raggedright
Failure exposed or controlled
\end{minipage} & \begin{minipage}[b]{\linewidth}\raggedright
Evidence status
\end{minipage} \\
\midrule\noalign{}
\endhead
\bottomrule\noalign{}
\endlastfoot
Single-manager finality & One authority boundary for closure and durable mutation & split-brain doctrine, people, or closure changes & mechanism evidence \\
Role families & Repeated seats share interface contracts & bespoke prompts with incompatible handoff expectations & mechanism evidence \\
Pods or lanes & Upstream, review, execution, and replay paths are routable before work starts & every worker infers the next collaborator from context & mechanism evidence \\
Default downstream chains & Normal handoffs are structural, exceptions are recorded & hidden local reroutes and orphaned ownership & mechanism evidence \\
Capacity, overflow, and failover & Overload or unavailable owners become explicit routing conditions & silent accumulation of stalled AI work & mechanism evidence; bounded counts not reported \\
Production/research separation & Method learning cannot bypass live authority & research output becomes unauthorized production action & mechanism evidence \\
Truth-surface synchronization & Role text, profile, tool set, work record, playbook, and activation agree before a change is live & a worker exists on paper but cannot operate correctly & mechanism evidence \\
\end{longtable}
}

\textbf{Limitation.} These are topology mechanisms observed in one organization; the report does not measure topology performance or claim a generally optimal structure.

\subsection{8. Anchor Mechanism Cases}\label{anchor-mechanism-cases}

The case studies use sanitized labels. They are organized to show the main institutional-control mechanisms across state, authority, communication, and learning, not private operational details.

The main body keeps five anchor cases that cover the paper's central mechanisms; additional cases are retained in Appendix J for completeness.

{\def\LTcaptype{none} % do not increment counter
\begin{longtable}[]{@{}
  >{\raggedright\arraybackslash}p{(\linewidth - 10\tabcolsep) * \real{0.1667}}
  >{\raggedright\arraybackslash}p{(\linewidth - 10\tabcolsep) * \real{0.1667}}
  >{\raggedright\arraybackslash}p{(\linewidth - 10\tabcolsep) * \real{0.1667}}
  >{\raggedright\arraybackslash}p{(\linewidth - 10\tabcolsep) * \real{0.1667}}
  >{\raggedright\arraybackslash}p{(\linewidth - 10\tabcolsep) * \real{0.1667}}
  >{\raggedright\arraybackslash}p{(\linewidth - 10\tabcolsep) * \real{0.1667}}@{}}
\toprule\noalign{}
\begin{minipage}[b]{\linewidth}\raggedright
Case
\end{minipage} & \begin{minipage}[b]{\linewidth}\raggedright
Location
\end{minipage} & \begin{minipage}[b]{\linewidth}\raggedright
Mechanism
\end{minipage} & \begin{minipage}[b]{\linewidth}\raggedright
Claim type
\end{minipage} & \begin{minipage}[b]{\linewidth}\raggedright
Evidence basis
\end{minipage} & \begin{minipage}[b]{\linewidth}\raggedright
Limitation
\end{minipage} \\
\midrule\noalign{}
\endhead
\bottomrule\noalign{}
\endlastfoot
A & main & manager closure for complete-but-open work & bounded metric & S1 closeout audit & no full before/after rate \\
B & Appendix J & message freshness and no-reply suppression & extracted mechanism & M1 packets, S1 closeouts & not inbox census \\
C & main & verifier-reconstructible artifact and action gate & extracted mechanism & verifier comments & domain fields abstracted \\
D & Appendix J & targeted doctrine patching & mechanism evidence & sanitized patch pattern & no cost, error, or recurrence claim \\
E & main & replay-to-mutation classification & extracted mechanism & replay and retro records & no recurrence rate \\
F & Appendix J & adaptive topology, role, tool, and staffing change & extracted mechanism & topology cases & no optimum claim \\
G & Appendix J & replay without immediate durable mutation & negative mechanism & S1 retro closeout & no immediate mutation \\
H & main & sprint compilation into child-path work & extracted mechanism & sprint structure & selected windows only \\
I & Appendix J & role-family capacity expansion with activation gating & extracted mechanism & capacity case & no recurrence or safety claim \\
J & main & durable artifact growth from repeated local fixes & mechanism synthesis & Cases A-I, Appendix I intervention timeline, and runtime/tool-surface evidence & no recurrence claim \\
\end{longtable}
}

\subsubsection{8.1 Case A: Manager Verification Queue And Atomic Closeout}\label{case-a-manager-verification-queue-and-atomic-closeout}

\textbf{Initial failure.} Work appeared substantively complete, but the governing record did not reach a valid final state. Some items were complete-but-open; others were review-ready but still mixed with active work. The organization could tell a story of completion, but the workboard did not carry closure truth.

\textbf{Organization diagnosis.} A worker-level view would ask whether the assigned output was produced. The organizational failure was invalid closure: no authorized owner had accepted the final state, checked child/replay obligations, and left no ambiguous next action.

\textbf{Intervention.} The organization introduced manager verification queue discipline. A manager review cycle had to end in exactly one state: closed, kept open for named dependency, rejected with missing piece, or rerouted to the exact next owner.

\textbf{Changed durable layer.} Workboard process, manager role, closure doctrine, and runtime-contract rules.

\textbf{Evidence artifact.} Extracted work-record cases show complete-but-open and umbrella-closeout failures, plus follow-through where remaining cases were closed or left open for named dependencies.

\textbf{General principle.} Completion is not closure. AI organizations benefit from closure validity as a first-class institutional state.

\textbf{Limitation.} Available evidence supports the mechanism but does not include a full before/after closure-validity rate.

\subsubsection{8.2 Case C: From Plausible Prose To Verifier-Reconstructible Artifact And Action Gate}\label{case-c-from-plausible-prose-to-verifier-reconstructible-artifact-and-action-gate}

\textbf{Initial failure.} A domain specialist could produce plausible narrative analysis, but downstream review needed reconstructible fields: source basis, freshness, settlement mapping, uncertainty, probability or decision structure, confidence, contradiction checks, and no-action triggers.

\textbf{Organization diagnosis.} Fluency hid two missing institutional gates: evidence sufficiency for review and action authorization for downstream execution.

\textbf{Intervention.} Research outputs became structured packets. Review playbooks defined packet quality gates and action gates. Verifier messages asked for stateful decisions: approve, repair, reroute, decline, or request manager review.

\textbf{Changed durable layer.} Playbook doctrine, verifier role, work-record acceptance criteria.

\textbf{Evidence artifact.} Extracted work records and verifier inbox samples show packet-gate failures, repair routing, approvals that modified the downstream envelope, and declines that explicitly withheld execution authorization. The case shows two gates: whether the artifact was reconstructible and whether downstream action was authorized.

\textbf{General principle.} Technical AI work benefits from verifier-reconstructible artifacts, not persuasive prose, and from action gates that treat a good artifact as insufficient for execution authority.

\textbf{Limitation.} Domain-specific packet fields are abstracted for readers outside trading.

\subsubsection{8.3 Case E: Replay-To-Institutional-Mutation}\label{case-e-replay-to-institutional-mutation}

\textbf{Initial failure.} Completed work could produce outcome truth later, but the lesson could remain trapped in a ticket, message, or one worker's context. Conversely, one-off outcome facts could be over-promoted into general doctrine.

\textbf{Organization diagnosis.} Reflection and memory can generate lessons, but they do not decide where the organization should change.

\textbf{Intervention.} Replay records reconstructed the original path and classified learning destination: role, playbook, tool, staffing, topology, workboard, message, or case-only. The stronger replay form also records the role and playbook versions active at decision time, so expected-versus-realized review does not judge old work only by later doctrine. Durable changes required manager-employee authority inside the human-owner governance boundary.

\textbf{Changed durable layer.} Replay doctrine, role/playbook/tool/topology mutation process, manager authority boundary.

\textbf{Evidence artifact.} Extracted replay and retro cases show scoped learning, case-only decisions, and selected durable-change mechanisms, while S1 shows a negative case where replay did not immediately become durable mutation.

\textbf{General principle.} Learning is incomplete until the durable fix lands in the correct organizational layer.

\textbf{Limitation.} Recurrence after mutation is only partially extracted. This report does not claim that any replay-driven change removed a failure class.

\subsubsection{8.4 Case H: Sprint Compilation Into Child Path Work}\label{case-h-sprint-compilation-into-child-path-work}

\textbf{Initial failure.} A broad owner request could sound operationally clear while leaving the actual AI labor frontier underspecified: which paths existed, who owned each path, which gate applied, which employee was activated, which paths required replay, and when the parent could close.

\textbf{Organization diagnosis.} A workflow view could treat the sprint as a set of tasks. The organizational failure was that owner intent had not been compiled into durable state that every employee could route, verify, close, or replay without relying on private context.

\textbf{Intervention.} The manager compiled broad intent into a live umbrella, retro umbrella, child path records, shared control fields, direct activation messages, verifier gates, no-action or closeout states, and replay obligations.

\begin{Shaded}
\begin{Highlighting}[]
\NormalTok{owner intent}
\NormalTok{  {-}\textgreater{} live umbrella}
\NormalTok{  {-}\textgreater{} retro umbrella}
\NormalTok{  {-}\textgreater{} child path records}
\NormalTok{  {-}\textgreater{} direct activations}
\NormalTok{  {-}\textgreater{} verifier gates}
\NormalTok{  {-}\textgreater{} close / no{-}action / replay}
\end{Highlighting}
\end{Shaded}

\textbf{Changed durable layer.} Workboard topology, message activation, manager closeout, and replay surface.

\textbf{Evidence artifact.} S1/R1/V1 sprint windows show live and retro pairings, child path records, manager-reviewed final dispositions, and ordered-path replay separation. The broader P1/B1 spot-check is used only as context, not as a full workboard-frontier rate.

\textbf{General principle.} The workboard is not a passive tracker. It is the institutional state machine that makes vague intent executable by bounded AI employees.

\textbf{Limitation.} This report shows this mechanism in selected and spot-check samples, but it does not prove the compilation topology is optimal or sufficient for all domains.

\subsubsection{8.5 Case J: Durable Artifact Growth From Repeated Failure}\label{case-j-durable-artifact-growth-from-repeated-failure}

\textbf{Initial failure.} Repeated work initially relied on bespoke prompts, comments, or case-specific work text. A manager or employee could correct a local problem, but the same missing structure could reappear because the correction had not become a durable institutional artifact.

\textbf{Organization diagnosis.} The failure was not only weak instruction following. It was that one-off correction was not durable doctrine, authority, state, or activation protocol. The organization needed a way to decide which repeated lessons belonged in role contracts, work-record fields, message rules, verifier decision forms, playbooks, replay doctrine, or tool surfaces.

\textbf{Intervention.} Repeated fixes were promoted into durable artifact forms by the manager employee, inside a human-owner governance boundary. Repeated role ambiguity became role-family contracts. Repeated case-truth ambiguity became work-record fields. Missed handoffs and false motion became message rules. Plausible but unauthorized artifacts became verifier decision forms. Outcome-learning gaps became replay doctrine. Capability gaps became tool-surface changes.

\textbf{Changed durable layer.} Role contracts, workboard/work-record structure, message protocol, verifier gate fields, playbook doctrine, replay process, and tool surface.

\textbf{Evidence artifact.} The intervention timeline, replay cases, targeted doctrine patching, workboard compilation, message discipline, verifier gates, and role-family activation case all show local failures being routed toward durable artifact changes rather than left as isolated conversation.

\textbf{General principle.} Durable artifact forms are institutional learning outputs. The contribution is not that complete schemas existed at the start, but that runtime-bounded work made missing structure visible enough for authorized promotion.

\textbf{Limitation.} This is mechanism evidence. It does not show unconstrained self-modification or recurrence reduction.

\subsection{9. Discussion}\label{discussion}

\subsubsection{9.1 Main Interpretation}\label{main-interpretation}

The main interpretation is that this field study suggests high-skill AI labor may need a control layer around agents. The organization did not establish global decision-quality improvement merely because workers received more instructions. It made reliability-relevant failures more visible and controllable when failures were converted into durable primitives: role authority, playbook doctrine, work-state truth, message lifecycle, evidence gates, no-action states, replay, and authorized mutation.

State, authority, communication, and learning are not claimed to be a complete taxonomy of AI organizations. They are the four control surfaces that became visible in this field study and that organized the observed failure-to-intervention pattern.

This is why ``AI organization'' should not be used as a loose metaphor for a group of agents. It should name a concrete institutional architecture.

The Related Work positioning table summarizes why agent SDKs, MCP, A2A, enterprise platforms, workflow graphs, and benchmarks remain useful but incomplete for this institutional-control question.

The deployment was AI-staffed for daily operational labor and human-owned for goals, risk boundaries, and outer governance. The manager employee performed sprint, closeout, replay, and institutional-mutation work within that boundary. In the S1 window, the bounded sample includes one sprint-opening owner directive, 20 manager-reviewed final closeouts, at least 9/20 child paths requiring packet/source repair before final decision, and no direct staffing/tool/topology mutation. This is evidence about a human-owned, AI-staffed institution, not a claim of fully unsupervised replacement.

\subsubsection{9.2 Falsifiers And Boundary Conditions}\label{falsifiers-and-boundary-conditions}

The thesis would weaken under evidence that:

\begin{itemize}
\tightlist
\item
  improvements came mostly from model upgrades rather than organizational interventions;
\item
  human intervention remained the hidden driver of correctness;
\item
  work records became compliant but decisions did not improve;
\item
  no-action states hid indecision rather than valid restraint;
\item
  replay did not reduce recurrence or improve doctrine;
\item
  the mechanisms failed outside prediction-market trading;
\item
  manager finality became an unscalable bottleneck.
\end{itemize}

The paper should welcome these tests. Institutional control is an empirical claim, not a slogan.

\subsubsection{9.3 Practical Design Rules For AI Organizations}\label{practical-design-rules-for-ai-organizations}

The practical design rules are:

\begin{itemize}
\tightlist
\item
  The model is not the system of record.
\item
  Tool access is not authority.
\item
  Messages are work triggers, not chat.
\item
  Review must produce a state decision, not commentary.
\item
  No-action must be a valid audited output.
\item
  Learning is incomplete until it changes the right durable layer or is classified case-only.
\item
  Human owners move upward into goals, risk boundaries, and outer governance; manager employees perform scoped institutional operation and mutation inside those boundaries.
\end{itemize}

\subsubsection{9.4 Humanlike Interface, Machine-Native Control}\label{humanlike-interface-machine-native-control}

The paper uses terms such as employee, manager, sprint, review, handoff, and retro because they are legible handles for humans and for other AI workers. Those terms should not be read as anthropomorphic evidence or as roleplay. They are the interface layer.

The control layer is machine-native: authority order, role boundaries, tool surfaces, work-record state, child-frontier topology, message freshness, verifier gates, no-action states, fixed stamps, replay linkage, and manager-employee-authorized mutation. A title does not grant authority unless role text, tool access, governing record state, and activation all support it. A manager review is not a courtesy comment; it is a state transition. A sprint kickoff is not a social ritual; it is compiled into live and retro umbrellas, child paths, gates, and exact activations.

This interface/control separation is a useful way to read the whole system. The organization can look familiar, but its auditability comes from explicit state machinery.

\subsubsection{9.5 Prototype System As Reference Implementation}\label{prototype-system-as-reference-implementation}

The prototype system used in this study is best treated as one reference implementation of broader organization-protocol semantics. The paper does not imply that others must adopt the same product, implementation substrate, folder structure, or internal implementation.

The portable primitives are:

\begin{itemize}
\tightlist
\item
  employee principal;
\item
  role contract;
\item
  governing work record;
\item
  work message;
\item
  playbook doctrine;
\item
  verifier gate;
\item
  no-action state;
\item
  replay record;
\item
  doctrine patch;
\item
  institutional mutation.
\end{itemize}

The standardization opportunity is analogous to but above MCP. MCP gives AI applications a standard way to access tools, resources, and prompts. An organization protocol would give AI employees a standard way to operate under identity, authority, work state, doctrine, verification, replay, and mutation semantics.

A product-neutral object model would make those semantics explicit rather than binding them to one implementation:

Product-neutral means technology-independent: the object model is not tied to this prototype, any issue tracker, a trading workflow, or MCP. It is a vocabulary for institutional control objects and lifecycle rules that another implementation could store in a database, workflow engine, issue tracker, or tool server while preserving the same authority semantics.

{\def\LTcaptype{none} % do not increment counter
\begin{longtable}[]{@{}
  >{\raggedright\arraybackslash}p{(\linewidth - 6\tabcolsep) * \real{0.2500}}
  >{\raggedright\arraybackslash}p{(\linewidth - 6\tabcolsep) * \real{0.2500}}
  >{\raggedright\arraybackslash}p{(\linewidth - 6\tabcolsep) * \real{0.2500}}
  >{\raggedright\arraybackslash}p{(\linewidth - 6\tabcolsep) * \real{0.2500}}@{}}
\toprule\noalign{}
\begin{minipage}[b]{\linewidth}\raggedright
Object
\end{minipage} & \begin{minipage}[b]{\linewidth}\raggedright
Key fields
\end{minipage} & \begin{minipage}[b]{\linewidth}\raggedright
Lifecycle
\end{minipage} & \begin{minipage}[b]{\linewidth}\raggedright
Why tool protocols alone do not cover it
\end{minipage} \\
\midrule\noalign{}
\endhead
\bottomrule\noalign{}
\endlastfoot
Organization & owner boundary, member principals, role families, work surfaces, doctrine surfaces, verifier topology & created, configured, staffed, mutated, audited & tool protocols expose tools and resources, not institutional membership, authority, finality, or doctrine ownership \\
Employee principal & identity, role, tools, inbox, status, history & hired, active, updated, suspended, terminated & tool protocols do not define persistent labor identity \\
Role contract & owned decisions, forbidden decisions, outputs, escalation, handoff interface & created, patched, versioned, retired & tool access is not authority \\
Playbook doctrine & scope, inputs, outputs, gates, no-action criteria, version & created, read, patched, versioned & retrieved documentation is not durable operating doctrine \\
Governing work record & owner, state, evidence, blocker, next action, verifier state & opened, routed, blocked, closed, replayed & tool protocols do not define case truth or closure \\
Work message & sender, recipient, work item, ask, reply flag, freshness state & sent, actionable, superseded, already done & chat is not work-state transfer \\
Verifier gate & artifact status, criteria, decision, execution authorization, next owner, authority effect & accept, repair, reroute, decline, needs review & validation alone is not institutional authority \\
No-action state & type, explicit zero, evidence basis, authority basis, next owner/action, replay requirement & proposed, verified, closed, replayed, classified case-only & capability protocols do not define when not acting is valid work \\
Closure decision & accepting authority, final state, dependency, rejection/reroute reason, replay obligation & proposed, accepted, kept open, rejected, rerouted, closed & execution protocols do not define institutional finality \\
Replay record & original path, outcome truth, role/playbook versions, learning destination & pending, completed, mutated, classified case-only & traces do not decide doctrine change \\
Institutional mutation & trigger, source evidence, destination layer, authorized actor, mutation type, activation step, expected behavior change, recurrence check & proposed, authorized, applied, activated, rechecked & administrative APIs do not define governed adaptation \\
Doctrine patch & target artifact, old string, new string, reason, authority, version result & proposed, applied, rejected, versioned & text editing is not enough for auditable institutional mutation \\
\end{longtable}
}

For example, in a synthetic security-incident desk, an \texttt{employee\ principal} can be a persistent triage analyst, a \texttt{role\ contract} can permit alert classification while forbidding containment approval, a \texttt{governing\ work\ record} can own incident truth, a \texttt{verifier\ gate} can accept a packet as reconstructible while declining containment authorization, a \texttt{no-action\ state} can record monitor-only handling, a \texttt{closure\ decision} can close the triage phase while replay stays open, and an \texttt{institutional\ mutation} can update a playbook or role boundary after repeated evidence. The public synthetic example keeps this walkthrough at the organization-protocol level rather than giving security-operations instructions.

\subsubsection{9.6 Why The Prediction-Market Desk Is A Useful Critical Case}\label{why-the-prediction-market-desk-is-a-useful-critical-case}

Prediction-market trading is a strong empirical stress test because it combines uncertainty, external evidence, delayed outcomes, and real action. The domain punishes unsupported action and makes no-action meaningful.

The general mechanisms transfer to other domains with:

\begin{itemize}
\tightlist
\item
  staged technical work;
\item
  external evidence;
\item
  role authority;
\item
  review and acceptance criteria;
\item
  delayed outcome truth;
\item
  operational consequences;
\item
  repeated workflow learning.
\end{itemize}

Examples include software maintenance, security operations, research operations, legal and compliance workflows, healthcare administration, financial research, and supply-chain incident response.

\subsection{10. Ethics, Risk, And Compliance}\label{ethics-risk-and-compliance}

This deployment operated in a real-money prediction-market setting. The paper therefore avoids presenting the work as trading advice, strategy disclosure, or evidence that AI systems may act without institutional controls.

This paper is not financial advice, not a trading strategy, and not a recommendation to deploy autonomous trading agents.

The paper intentionally omits market names, order sizes, account details, strategy thresholds, and timing details that could expose operational strategy or encourage imitation.

The relevant risk controls are:

\begin{itemize}
\tightlist
\item
  human owner boundary;
\item
  role-based execution authority;
\item
  verifier gates before action;
\item
  no-action states;
\item
  risk and exposure constraints;
\item
  manager closure;
\item
  replay and auditability;
\item
  separation between live closure and delayed outcome truth;
\item
  privacy-preserving evidence extraction.
\end{itemize}

The ethical claim is not that autonomy is inherently desirable. It is that if AI systems are used for high-skill labor, the responsible path is to make authority, evidence, review, stopping, and learning explicit.

\subsection{11. Limitations And Threats To Validity}\label{limitations-and-threats-to-validity}

\subsubsection{11.1 Internal Validity}\label{internal-validity}

The organization changed while it operated. This is realistic for applied deployment but limits clean causal inference. Interventions overlap, and this study does not estimate exact before/after windows. The system owner also shaped the institution, selected the research framing, and directed parts of the evidence extraction process. This is normal for an applied field study, but it means the paper should avoid claiming neutral observation from outside the system.

\subsubsection{11.2 External Validity}\label{external-validity}

The empirical setting is one AI organization in one high-skill domain. The paper can support depth and mechanism discovery, not general causal certainty. Prediction markets provide unusually clear settlement and outcome truth; other domains may have slower, noisier, or more contested feedback. Single-manager finality also prevents split-brain mutation but may limit scale. Subsequent work could explore delegated managers, domain managers, scoped mutation authority, or quorum review.

\subsubsection{11.3 Measurement Validity}\label{measurement-validity}

The evidence emphasizes high-signal cases plus selected bounded windows. This study does not estimate the broader denominators required for organization-wide rate claims. Roles and playbooks remain natural-language artifacts; explicitness helps but does not guarantee compliance. Future systems may need machine-checkable doctrine or typed organizational policies.

Duplicate/canonical-parent repair is source-backed and supports the workboard graph-truth mechanism, but the available extraction is not sufficient for a bounded quantitative workboard-frontier result.

Stronger evidence would require a full-corpus extraction, independent coding, pre/post windows tied to specific interventions, and recurrence checks after activated doctrine or topology changes.

\subsubsection{11.4 Coding Reliability}\label{coding-reliability}

Evidence extraction and coding were author-directed. Subsequent work could use independent coders or blinded review for a subset of work records, messages, verifier decisions, no-action states, and replay records.

\subsubsection{11.5 Human-Owner Boundary}\label{human-owner-boundary}

The organization was AI-staffed in daily labor, but the human owner remained an authority boundary. The paper should not claim fully unsupervised human-free operation. The observed system is best understood as a human-owned, AI-staffed institution.

\subsubsection{11.6 Reproducibility}\label{reproducibility}

Privacy and operational sensitivity limit how much raw evidence can be published. Public artifacts should use sanitized field sets, aggregate metrics, and publication-ready case labels.

The public package lets readers inspect and try the institutional-control semantics without exposing private trading details. It includes the technical report with a product-neutral organization-protocol object model, API/MCP documentation for sandboxed use, and a synthetic non-trading workboard example. The report itself includes paper-safe artifact field sets for role contracts, work records, messages, verifier decisions, closure decisions, and replay records. These field sets are distilled from evolved institutional artifacts, not presented as fixed external templates. In the deployment, such forms matured through repeated institutional use, failure pressure, verifier feedback, and doctrine mutation. The public API demonstrates institutional semantics, not autonomous trading.

The synthetic example is a demonstration of organization-protocol semantics, not empirical evidence and not a benchmark.

\subsection{12. Conclusion}\label{conclusion}

The central lesson of this deployment is that agent capability is not the same as auditable labor control. This study suggests that high-skill AI work may benefit from durable institutional control: employees, authority, work records, messages, doctrine, evidence gates, no-action states, replay, and governed adaptation.

The AI-staffed prediction-market desk began as a practical experiment in AI labor and became evidence for a broader claim. Many failures were not addressed by asking agents to think harder. They were addressed by changing the organization around the agents.

The next frontier is not merely agents that can act. It is organizations that can make AI action accountable.

\subsection{13. Figures, Tables, And Appendices}\label{figures-tables-and-appendices}

\subsubsection{13.1 Included Figures}\label{included-figures}

Main-body figures included in this report:

\begin{enumerate}
\def\labelenumi{\arabic{enumi}.}
\tightlist
\item
  \textbf{Figure 1: From Models To AI Organizations.}
\item
  \textbf{Figure 2: Institutional Control Stack.}
\item
  \textbf{Figure 3: Workboard As Institutional State Machine.}
\item
  \textbf{Figure 4: Creator -\textgreater{} Verifier -\textgreater{} Executor -\textgreater{} Manager -\textgreater{} Replay.}
\item
  \textbf{Figure 5: Work Item Lifecycle.}
\end{enumerate}

The report focuses on five main figures to keep the institutional-control argument compact.

\subsubsection{13.2 Appendix A: Role Contract Field Set}\label{appendix-a-role-contract-field-set}

Paper-safe field set distilled from evolved institutional artifacts:

\begin{itemize}
\tightlist
\item
  mandate;
\item
  scope;
\item
  owned decisions;
\item
  forbidden decisions;
\item
  required outputs;
\item
  upstream inputs;
\item
  downstream handoffs;
\item
  escalation triggers;
\item
  non-goals.
\end{itemize}

\subsubsection{13.3 Appendix B: Work Record Field Set}\label{appendix-b-work-record-field-set}

Paper-safe field set distilled from evolved institutional artifacts:

\begin{itemize}
\tightlist
\item
  work item label;
\item
  owner role;
\item
  status;
\item
  parent umbrella or child path label;
\item
  gate state;
\item
  evidence summary;
\item
  verifier state;
\item
  blocker;
\item
  next owner;
\item
  next action;
\item
  replay state;
\item
  closure basis.
\end{itemize}

\subsubsection{13.4 Appendix C: Work Message Field Set}\label{appendix-c-work-message-field-set}

Paper-safe field set distilled from evolved institutional artifacts:

\begin{itemize}
\tightlist
\item
  sender role;
\item
  recipient role;
\item
  work item;
\item
  exact ask;
\item
  exact next action;
\item
  reply needed: yes/no;
\item
  freshness state;
\item
  outcome;
\item
  messages suppressed or sent.
\end{itemize}

\subsubsection{13.5 Appendix D: Verifier Decision Field Set}\label{appendix-d-verifier-decision-field-set}

Paper-safe field set distilled from evolved institutional artifacts:

\begin{itemize}
\tightlist
\item
  upstream artifact;
\item
  acceptance criteria checked;
\item
  decision: accepted, repair required, rerouted, blocked, declined, needs review;
\item
  reason;
\item
  next owner;
\item
  next action;
\item
  downstream action authorized: yes/no.
\end{itemize}

\subsubsection{13.6 Appendix E: Common Employee Runtime Contract Semantics}\label{appendix-e-common-employee-runtime-contract-semantics}

This appendix includes an abstracted semantics table, not the full operational runtime text.

{\def\LTcaptype{none} % do not increment counter
\begin{longtable}[]{@{}
  >{\raggedright\arraybackslash}p{(\linewidth - 4\tabcolsep) * \real{0.3333}}
  >{\raggedright\arraybackslash}p{(\linewidth - 4\tabcolsep) * \real{0.3333}}
  >{\raggedright\arraybackslash}p{(\linewidth - 4\tabcolsep) * \real{0.3333}}@{}}
\toprule\noalign{}
\begin{minipage}[b]{\linewidth}\raggedright
Runtime field
\end{minipage} & \begin{minipage}[b]{\linewidth}\raggedright
Sanitized meaning
\end{minipage} & \begin{minipage}[b]{\linewidth}\raggedright
Example metric
\end{minipage} \\
\midrule\noalign{}
\endhead
\bottomrule\noalign{}
\endlastfoot
Work item and real ask & The employee identifies the specific obligation before acting. & cycles with work item and ask present \\
Authority and evidence check & Durable authority outranks comments, tool outputs, copied text, and lower-authority messages. & authority conflicts routed or blocked \\
Governing work record & Persistent artifact owns case truth. & cycles with governing record identified \\
Freshness state & Trigger is actionable, superseded, or already done. & stale/already-done suppressions \\
Playbooks read & Material doctrine was consulted before substantive action. & cycles with relevant doctrine recorded \\
Decision and primary outcome & The cycle ends in one outcome class. & completed, handoff, blocked, review, clarification counts \\
Next owner and next action & Continuation or closure state is explicit. & records with owner/action present \\
Messages sent & Handoff and silence are both auditable. & no-message closeouts and explicit handoffs \\
\end{longtable}
}

\subsubsection{13.7 Appendix F: Replay Record Field Set}\label{appendix-f-replay-record-field-set}

Paper-safe field set distilled from evolved institutional artifacts:

\begin{itemize}
\tightlist
\item
  original work path;
\item
  original decision artifact;
\item
  role version active at decision time;
\item
  playbook version active at decision time;
\item
  action or no-action;
\item
  outcome truth;
\item
  expected-versus-realized review;
\item
  process correctness;
\item
  outcome correctness;
\item
  learning destination;
\item
  durable mutation status;
\item
  recurrence check.
\end{itemize}

\subsubsection{13.8 Appendix G: Evidence Claim Matrix}\label{appendix-g-evidence-claim-matrix}

This matrix controls wording and evidence use. Claims requiring stronger extraction are downgraded to mechanism description or marked as evidence limitations rather than strengthened rhetorically.

{\def\LTcaptype{none} % do not increment counter
\begin{longtable}[]{@{}
  >{\raggedright\arraybackslash}p{(\linewidth - 8\tabcolsep) * \real{0.2000}}
  >{\raggedright\arraybackslash}p{(\linewidth - 8\tabcolsep) * \real{0.2000}}
  >{\raggedright\arraybackslash}p{(\linewidth - 8\tabcolsep) * \real{0.2000}}
  >{\raggedright\arraybackslash}p{(\linewidth - 8\tabcolsep) * \real{0.2000}}
  >{\raggedright\arraybackslash}p{(\linewidth - 8\tabcolsep) * \real{0.2000}}@{}}
\toprule\noalign{}
\begin{minipage}[b]{\linewidth}\raggedright
Claim
\end{minipage} & \begin{minipage}[b]{\linewidth}\raggedright
Allowed wording
\end{minipage} & \begin{minipage}[b]{\linewidth}\raggedright
Evidence support
\end{minipage} & \begin{minipage}[b]{\linewidth}\raggedright
Evidence label
\end{minipage} & \begin{minipage}[b]{\linewidth}\raggedright
Forbidden wording
\end{minipage} \\
\midrule\noalign{}
\endhead
\bottomrule\noalign{}
\endlastfoot
Four-pillar institutional control & State, authority, communication, and learning are a useful organizing frame for the observed mechanisms. & Framework synthesis plus Results organization. & Conceptual claim & These four pillars are complete, universal, or proven sufficient. \\
State control & Made work truth more auditable in S1. & S1 closeout records and fixed stamps. & Bounded metric claim & Established a global reliability rate. \\
Single-step scope control & Each employee cycle is bounded to one authorized outcome, reducing scope ambiguity and making audit easier. & Harness evidence and S1 final closeout records. & Mechanism evidence claim & Proved over-expansion reduction, reliability improvement, or complete scope control. \\
Authority control & Turned review into state decisions in observed verifier paths. & S1 verifier comments and reroutes. & Bounded metric claim & Established complete autonomous execution control. \\
Action-gate decline & Verifier topology separates artifact review from execution authorization in observed verifier paths. & Verifier comments and no-execution dispositions. & Mechanism evidence claim & Verifier decline proves safety or improved decision quality. \\
Workboard state machine & Compiled owner intent into live/retro umbrellas, child paths, activations, gates, and replay obligations. & S1/R1/V1 sprint structure and workboard evolution. & Mechanism evidence claim & Proved the topology is optimal or universally required. \\
Durable artifact growth & Recurring fixes were promoted into durable artifact forms by the manager employee inside a human-owner governance boundary. & Cases A-J plus runtime/tool-surface evidence. & Mechanism evidence claim & The system autonomously bootstrapped itself without manager-employee authority, proved self-improvement, or reduced recurrence. \\
False-motion control & Made stale, duplicate, and non-material work suppressible in observed packets and closeouts. & M1 packets and S1 closeouts. & Mechanism evidence claim & Showed there are no remaining loops. \\
Learning control & Separated stopping, replay, and doctrine-change destinations. & S1 replay records and retro closeout. & Bounded metric claim & Established recurrence reduction. \\
Explicit zero & Explicit zero is a verified no-action state where the organization records that no live path or executable action remains. & No-action mechanism evidence and workboard state evidence. & Mechanism evidence claim & Explicit zero proves correctness or successful outcome. \\
Topology control & Made routing and authority more explicit through role families, lanes, default chains, manager finality, and synchronized truth surfaces. & Organization evolution and mechanism cases. & Mechanism evidence claim & Established a general best organizational topology. \\
Person-role-aware tool surface & Shows tool capability interpreted through employee principal, role, work state, and manager-employee / human-owner authority boundaries. & Tool-surface evolution and authority-control cases. & Mechanism evidence claim & Claimed tool protocols alone enforce all institutional authority. \\
Institutional control & Suggested as useful for high-skill AI labor. & Field study plus mechanism evidence. & Conceptual claim & Claimed this is mandatory for all reliable AI labor. \\
Durable mutation & Describes a mechanism for changing doctrine, roles, tools, staffing, topology, or message protocol. & Selected evolution/tool records. & Mechanism evidence claim & Reports aggregate mutation rates. \\
Authorized institutional mutation / role-family activation gating & Shows one source-backed role-family capacity mechanism where activation waited until required role, tool, routing, and activation surfaces were directly evidenced as aligned. & Case I, role-family capacity expansion with activation gating. & Mechanism evidence claim & Recurrence reduction, causal reliability improvement, complete activation safety, safe autonomous rollout, successful rollout. \\
Targeted doctrine patching & Identifies anchored edits as a mechanism for local doctrine change. & Source-backed tool behavior and sanitized mechanism example. & Mechanism evidence claim; bounded quantitative effects not estimated & Quantifies context-cost savings or error reduction. \\
Reliability & Made reliability-relevant failures more visible and controllable. & State, authority, communication, and learning evidence. & Design implication & Established global decision-quality improvement. \\
\end{longtable}
}

\subsubsection{13.9 Appendix H: Evidence Status And Coding Guide}\label{appendix-h-evidence-status-and-coding-guide}

Current evidence is strongest for mechanism tracing and bounded mature-state audits. It is weaker for aggregate rates.

{\def\LTcaptype{none} % do not increment counter
\begin{longtable}[]{@{}
  >{\raggedright\arraybackslash}p{(\linewidth - 4\tabcolsep) * \real{0.3333}}
  >{\raggedright\arraybackslash}p{(\linewidth - 4\tabcolsep) * \real{0.3333}}
  >{\raggedright\arraybackslash}p{(\linewidth - 4\tabcolsep) * \real{0.3333}}@{}}
\toprule\noalign{}
\begin{minipage}[b]{\linewidth}\raggedright
Evidence type
\end{minipage} & \begin{minipage}[b]{\linewidth}\raggedright
Current status
\end{minipage} & \begin{minipage}[b]{\linewidth}\raggedright
Use in this report
\end{minipage} \\
\midrule\noalign{}
\endhead
\bottomrule\noalign{}
\endlastfoot
Mechanism case packets from work records & extracted & main qualitative results \\
Message evidence packets & extracted & message and verifier cases \\
S1 direct activation coverage & extracted mechanism evidence & 20/20 child records direct-activated from S1 live umbrella; not a full inbox census \\
Single-step scope control & mechanism evidence extracted & runtime-contract rule plus S1 final closeout records; no over-expansion reduction rate \\
Intervention timeline & extracted/inferred mix & empirical backbone \\
Harness evolution & inferred from mature runtime contract plus linked mechanism evidence & method claim with caveat; historical prompt-version sequence was not fully reconstructed \\
Metric definitions & defined & evaluation framework \\
Denominator-bounded metric values & extracted for selected sprint windows & sample-labeled Results counts \\
Broader P1/B1 spot-check & shallow metadata coding & selection-bias context, not full-corpus rates \\
Case I, role-family capacity expansion with activation gating & extracted mechanism evidence & activation-gating mechanism; no recurrence, reliability, or safety rate \\
Case J, durable artifact growth & mechanism evidence extracted & recurring fixes promoted into durable artifact forms; no self-modification, recurrence, or generalization rate \\
Role-family activation readiness & source artifacts identified; bounded estimates not reported & includes staffing/role-family capacity, role contract readiness, work-tool assignment, routing, and direct activation dependencies \\
Recurrence windows & partial & this study does not estimate recurrence after mutation \\
Targeted doctrine patch quantitative evidence & defined but not estimated in this study & method claim with explicit caveat \\
\end{longtable}
}

Harness evolution is inferred from the mature runtime contract and linked mechanism evidence; the historical prompt-version sequence was not fully reconstructed.

For metrics defined but not estimated in this study, the coding schema is:

\begin{itemize}
\tightlist
\item
  current owner present: yes/no;
\item
  next action present: yes/no;
\item
  role authority respected: yes/no/unclear;
\item
  required evidence present: yes/no/partial;
\item
  verifier decision type: accept/repair/reroute/decline/block/needs review;
\item
  closure validity: valid/invalid/open/unclear;
\item
  no-action type: none/no-send/decline/stand-down/blocked/already-done/explicit zero;
\item
  replay state: not required/pending/complete/missing;
\item
  mutation destination: role/playbook/tool/staffing/topology/workboard/message/case-only/none;
\item
  recurrence: none observed/observed/not enough exposure.
\end{itemize}

\subsubsection{13.10 Appendix I: Intervention Timeline}\label{appendix-i-intervention-timeline}

The full intervention timeline records the durable layers that emerged as the organization moved from weak agent-team forms to institutional-control forms.

{\def\LTcaptype{none} % do not increment counter
\begin{longtable}[]{@{}
  >{\raggedright\arraybackslash}p{(\linewidth - 8\tabcolsep) * \real{0.2000}}
  >{\raggedright\arraybackslash}p{(\linewidth - 8\tabcolsep) * \real{0.2000}}
  >{\raggedright\arraybackslash}p{(\linewidth - 8\tabcolsep) * \real{0.2000}}
  >{\raggedright\arraybackslash}p{(\linewidth - 8\tabcolsep) * \real{0.2000}}
  >{\raggedright\arraybackslash}p{(\linewidth - 8\tabcolsep) * \real{0.2000}}@{}}
\toprule\noalign{}
\begin{minipage}[b]{\linewidth}\raggedright
Stage
\end{minipage} & \begin{minipage}[b]{\linewidth}\raggedright
Failure pressure
\end{minipage} & \begin{minipage}[b]{\linewidth}\raggedright
Intervention
\end{minipage} & \begin{minipage}[b]{\linewidth}\raggedright
Changed durable layer
\end{minipage} & \begin{minipage}[b]{\linewidth}\raggedright
Evidence status
\end{minipage} \\
\midrule\noalign{}
\endhead
\bottomrule\noalign{}
\endlastfoot
T0 & Useful specialist agents lacked stable ownership and closure & Persistent AI employees & employee principal, role, inbox & source-backed mechanism evidence \\
T1 & Roles blurred recommendation, approval, execution, closure, and doctrine editing & Role contracts & role text and role families & extracted \\
T2 & Lessons stayed in tickets, messages, or memory & Playbooks as doctrine & playbook versions & extracted \\
T3 & Conversation appeared to own truth & Governing work records & workboard state & extracted \\
T4 & Assignment or comments did not reliably wake employees & Direct work messages & message protocol & extracted \\
T5 & Employees could act on stale messages & Freshness pre-check and suppression & runtime contract, inbox tools, messages & extracted \\
T6 & AI workers could self-certify plausible artifacts & Verifier topology & roles, playbooks, gates & extracted \\
T7 & Complete work remained open or review-ready without closure & Manager verification queue & workboard, manager role, runtime contract & extracted \\
T8 & Agents over-produced action & No-action states & review gates and work states & extracted \\
T9 & Live work stayed fake-active while awaiting delayed outcome truth & Replay split & workboard and replay records & extracted \\
T10 & Reflection remained commentary & Replay-to-mutation loop & mutation destinations & extracted \\
T11 & Long doctrine rewrites risked compression and drift & Targeted doctrine patches & role and playbook tools & mechanism evidence; bounded counts not reported \\
T12 & Capability and authority were confused & Dual tool plane and activation readiness & tool registry and role authority & extracted \\
T13 & More employees increased routing ambiguity & Pods and role families & organization topology & extracted \\
T14 & Fixes landed in the wrong layer or were not activated & Adaptive institution loop & institutional mutation process & extracted \\
\end{longtable}
}

Replay-to-mutation destinations can include roles, playbooks, tools, staffing, topology, workboards, or messages; the table uses the compact label \texttt{mutation\ destinations}.

\subsubsection{13.11 Appendix J: Additional Mechanism Cases}\label{appendix-j-additional-mechanism-cases}

These additional mechanism cases support the full case map while keeping the main body focused on the anchor cases.

\paragraph{Case B: Message Trigger Discipline And Stale/No-Reply Suppression}\label{case-b-message-trigger-discipline-and-staleno-reply-suppression}

\textbf{Initial failure.} Assignments, comments, and tracker fields did not reliably wake the correct employee. Later, direct messages themselves created a new problem: old or duplicate messages could cause repeated review, acknowledgment loops, or no-change state churn.

\textbf{Organization diagnosis.} Multi-agent chat treats communication as collaboration. The organizational failure was message lifecycle ambiguity: communication needed to be a work-transfer protocol with lifecycle state.

\textbf{Intervention.} Work messages were required to include a work item, exact ask, exact next action, and reply-needed flag. Freshness checks, using scoped inbox context, classified triggers as actionable, superseded, or already done. Routine replies were suppressed when no recipient needed to act.

\textbf{Changed durable layer.} Message protocol, transparent inbox tools, runtime-contract planning, loop-prevention doctrine.

\textbf{Evidence artifact.} Message evidence shows manager review requests with no-reply expectation, closeouts without direct replies, self-triggers that stopped after no-change checks, and execution requests rerouted because current gate truth no longer authorized action.

\textbf{General principle.} A message is not chat. It is an auditable trigger for one focused work obligation. Correct silence can be part of accountable, controllable AI labor.

\textbf{Limitation.} Message metrics combine two evidence types: work-record routing counts from S1 and packet examples from a bounded inbox sample. They support the mechanism, but they are not a global organization-wide loop-rate estimate.

\paragraph{Case D: Targeted Doctrine Patch Instead Of Whole Rewrite}\label{case-d-targeted-doctrine-patch-instead-of-whole-rewrite}

\textbf{Initial failure.} As role contracts and playbooks became more explicit, they also became longer. Whole-document rewrites for small changes risked compression, omitted clauses, restyling, accidental scope changes, and higher context and review cost.

\textbf{Organization diagnosis.} Prompt iteration treated doctrine as a single mutable blob. The organizational problem was durable-law mutation: small changes needed to land without damaging unrelated clauses.

\textbf{Intervention.} The system introduced targeted replacement edits for role and playbook text. A small correction could specify the old string and the new string, with optional replace-all only when appropriate. This made doctrine mutation local, reviewable, and less likely to silently rewrite unrelated policy.

\textbf{Changed durable layer.} Role/playbook mutation tools and doctrine maintenance process.

\textbf{Evidence artifact.} Sanitized mechanism evidence identifies anchored replacement behavior. A local-patch pattern is:

\begin{Shaded}
\begin{Highlighting}[]
\NormalTok{Old clause: if another employee should act, update the work record.}
\NormalTok{New clause: if another employee must act, send a direct message naming the work item, exact ask, exact next action, and reply expectation.}
\NormalTok{Patch rule: apply only when the old clause is found exactly once.}
\end{Highlighting}
\end{Shaded}

An ambiguous-anchor pattern is:

\begin{Shaded}
\begin{Highlighting}[]
\NormalTok{Requested old clause: "send a direct message"}
\NormalTok{Result: rejected because the clause appeared in multiple doctrine sections.}
\NormalTok{Required repair: narrow the anchor or explicitly authorize replace{-}all.}
\end{Highlighting}
\end{Shaded}

\textbf{General principle.} AI organizations benefit from doctrine surgery, not only prompt rewriting.

\textbf{Limitation.} This case supports the mechanism but not quantitative claims about context-cost savings, lower edit error, or downstream recurrence reduction.

\paragraph{Case F: Adaptive Organization Through Role, Tool, Staffing, And Topology Change}\label{case-f-adaptive-organization-through-role-tool-staffing-and-topology-change}

\textbf{Initial failure.} Adding more AI workers increased routing ambiguity and tool/role mismatch. A worker could exist without the right institutional surface, or a pod could route work without the right verifier/executor chain.

\textbf{Organization diagnosis.} Adding specialists as prompts did not define staffing, authority, tool parity, activation, or routing topology.

\textbf{Intervention.} The organization developed role families, pods, default chains, manager-employee-authorized durable changes, tool parity checks, and activation discipline.

\textbf{Changed durable layer.} Organization topology, employee/tool lifecycle, role-family contracts, workboard routing.

\textbf{Evidence artifact.} Extracted organization rollout and activation-readiness cases show that staffing or tool repair could be the correct institutional fix, not another prompt revision.

\textbf{General principle.} An AI-staffed organization prototype can adapt institutionally: it can change people, doctrine, tools, routing, and topology.

\textbf{Limitation.} This case is presented as mechanism discovery from one organization, not evidence of a generally valid topology.

\paragraph{Case G: Replay Without Immediate Durable Mutation}\label{case-g-replay-without-immediate-durable-mutation}

\textbf{Initial failure.} A replay/retro closeout produced a research/doctrine follow-up artifact but did not immediately create a durable playbook, role, tool, staffing, topology, workboard, or message-protocol mutation.

\textbf{Organization diagnosis.} Institutional control made the learning gap visible, but visibility did not automatically mean the organization had completed durable learning. This is a negative case against overclaiming replay as recurrence reduction.

\textbf{Intervention.} Replay classification separated the case-level finding from immediate durable mutation.

\textbf{Changed durable layer.} Replay record and retro closeout state.

\textbf{Evidence artifact.} The S1 retro closeout recorded 0 immediate durable patches and 1 research/doctrine follow-up artifact.

\textbf{General principle.} Replay can separate outcome truth from live closure, but durable organizational learning remains incomplete until the mutation destination is resolved and the change is activated.

\textbf{Limitation.} This report does not measure whether the follow-up artifact later became durable doctrine or reduced recurrence.

\paragraph{Case I: Role-Family Capacity Expansion With Activation Gating}\label{case-i-role-family-capacity-expansion-with-activation-gating}

Case F describes the general topology pattern: role families, lanes, default chains, tool parity, and activation discipline. Case I shows one extracted role-family capacity lifecycle where activation was gated until the required surfaces aligned.

\textbf{Initial failure.} A time-sensitive work lane required additional same-role capacity, not just more prompts. The organization created or activated additional role-family seats only after the role contract, required work tools, routing path, and direct activation agreed. Until those surfaces aligned, the seat was not treated as live capacity.

\textbf{Organization diagnosis.} The activation boundary failed unless role, tool, routing, and activation surfaces all agreed.

\textbf{Intervention.} Activation was held until required role, tool, routing, and activation surfaces were directly evidenced as aligned.

\textbf{Changed durable layer.} Organization topology, role text, playbook/routing doctrine, tool/profile truth, activation protocol, and workboard cutover state.

\textbf{Evidence artifact.} Extracted mechanism evidence shows role-family capacity expansion with tool-class parity and activation-gating behavior. Sensitive tool details are omitted from public artifacts.

\textbf{General principle.} A new specialist or role-family seat is not live capacity until role authority, required work tools, routing, and direct activation agree.

\textbf{Limitation.} This case does not show recurrence effects, reliability effects, or a general rollout method. It shows one source-backed mechanism for holding activation until required role, tool, routing, and activation surfaces aligned.

\bibliography{references}

\end{document}
